\documentclass[11pt]{article}
\usepackage[margin=1in]{geometry}
\usepackage{hyperref}
\usepackage{graphicx}
\graphicspath{{./}}
\usepackage{natbib}
\usepackage{doi}
\setcitestyle{numbers,square} %
\makeatletter
\renewcommand{\hyper@natlinkbreak}[2]{#1}
\makeatother
\usepackage{titlesec}

\usepackage{amsmath}
\usepackage{amsthm}
\usepackage{dsfont} %
\theoremstyle{plain}
\newtheorem{theorem}{Theorem}
\newtheorem{proposition}[theorem]{Proposition}
\newtheorem{lemma}[theorem]{Lemma}
\newtheorem{corollary}[theorem]{Corollary}
\newtheorem{conjecture}[theorem]{Conjecture}

\theoremstyle{definition}
\newtheorem{definition}{Definition}
\newtheorem{axiom}[definition]{Axiom}
\newtheorem{assumption}[definition]{Assumption}
\newtheorem{criterion}[definition]{Criterion}
\newtheorem{example}[definition]{Example}

\theoremstyle{remark}
\newtheorem{remark}{Remark}
\newtheorem{observation}[remark]{Observation}
\newtheorem{property}[remark]{Property}

\usepackage{etoolbox}
\BeforeBeginEnvironment{theorem}{\needspace{4\baselineskip}\addvspace{6pt}}
\BeforeBeginEnvironment{corollary}{\needspace{4\baselineskip}\addvspace{6pt}}
\BeforeBeginEnvironment{lemma}{\needspace{4\baselineskip}\addvspace{6pt}}
\BeforeBeginEnvironment{proposition}{\needspace{4\baselineskip}\addvspace{6pt}}
\BeforeBeginEnvironment{definition}{\needspace{4\baselineskip}\addvspace{6pt}}
\BeforeBeginEnvironment{axiom}{\needspace{4\baselineskip}\addvspace{6pt}}
\BeforeBeginEnvironment{assumption}{\needspace{4\baselineskip}\addvspace{6pt}}
\BeforeBeginEnvironment{remark}{\needspace{4\baselineskip}\addvspace{9pt}}

\usepackage{longtable}
\usepackage{booktabs}
\usepackage{array}
\usepackage{calc}
\usepackage{caption}
\newenvironment{tablenote}
  {\par\addvspace{2pt}\footnotesize\noindent}
  {\par\addvspace{4pt}}
\newcounter{none}

\usepackage{parskip}

\widowpenalty=10000
\clubpenalty=10000
\usepackage[bottom]{footmisc}
\predisplaypenalty=10000
\postdisplaypenalty=150
\raggedbottom
\setlength{\skip\footins}{1.5\baselineskip}
\usepackage{needspace}
\let\origsection\section
\renewcommand{\section}{\needspace{12\baselineskip}\origsection}
\let\origsubsection\subsection
\renewcommand{\subsection}{\needspace{7\baselineskip}\origsubsection}
\let\origsubsubsection\subsubsection
\renewcommand{\subsubsection}{\needspace{4\baselineskip}\origsubsubsection}
\let\origparagraphcmd\paragraph
\renewcommand{\paragraph}{\needspace{3\baselineskip}\origparagraphcmd}
\let\origsubparagraph\subparagraph
\renewcommand{\subparagraph}{\needspace{3\baselineskip}\origsubparagraph}

\interdisplaylinepenalty=10000

\setcounter{secnumdepth}{3}

\titleformat{\paragraph}
  {\normalfont\normalsize\bfseries}{\theparagraph}{1em}{}
\titlespacing*{\paragraph}
  {0pt}{3.25ex plus 1ex minus .2ex}{1.5ex plus .2ex}
\titleformat{\subparagraph}
  {\normalfont\normalsize\bfseries}{\thesubparagraph}{1em}{}
\titlespacing*{\subparagraph}
  {0pt}{3.25ex plus 1ex minus .2ex}{1.5ex plus .2ex}

\providecommand{\tightlist}{%
  \setlength{\itemsep}{0pt}\setlength{\parskip}{0pt}}

\newcommand{\pandocbounded}[1]{%
  \sbox0{#1}%
  \ifdim\wd0>\linewidth
    \resizebox{\linewidth}{!}{#1}%
  \else
    #1%
  \fi
}

\usepackage{amssymb}
\usepackage{newunicodechar}
\newunicodechar{✓}{$\checkmark$}
\newunicodechar{≠}{$\neq$}
\newunicodechar{≡}{$\equiv$}
\newunicodechar{≈}{$\approx$}
\newunicodechar{δ}{$\delta$}
\newunicodechar{−}{$-$}
\newunicodechar{⁰}{$^{0}$}
\begingroup\renewcommand\PackageWarning[2]{}%
\newunicodechar{¹}{$^{1}$}
\newunicodechar{²}{$^{2}$}
\newunicodechar{³}{$^{3}$}
\endgroup
\newunicodechar{⁴}{$^{4}$}
\newunicodechar{⁵}{$^{5}$}
\newunicodechar{⁶}{$^{6}$}
\newunicodechar{⁷}{$^{7}$}
\newunicodechar{⁸}{$^{8}$}
\newunicodechar{⁹}{$^{9}$}
\newunicodechar{⁻}{$^{-}$}
\let\mystRightarrow\rightarrow
\renewcommand{\rightarrow}{\ensuremath{\mystRightarrow}}

\emergencystretch=2em

\hypersetup{colorlinks=true, allcolors=blue}

\title{Thermodynamics of Cooperation: Thermodynamic Friction}
\author{%
Keith Lostracco%
\\[0.25em]{\normalsize ORCID: \href{https://orcid.org/0009-0005-3106-006X}{0009-0005-3106-006X}}%
}
\date{}

\begin{document}
\maketitle

\begin{abstract}
We quantify the energy cost of adversarial resource competition between
self-maintaining open thermodynamic systems. Boundary constraints are
required for multi-agent resource allocation; we model what happens when
those constraints are violated. Using Tullock contest theory and a
boundary-integrity damage model, we prove that unconstrained conflict is
a net-negative thermodynamic outcome: under realistic physical
parameters, the total energy dissipated by conflict, including contest
expenditure, boundary damage, and thermodynamically irreversible repair
costs, exceeds the value of the contested resource. We derive explicit
Nash equilibrium investments and damage functions for both symmetric and
asymmetric contests, showing that contest dissipation approaches the
full resource value as the number of contestants grows. Extending the
analysis to cooperative networks, we prove that a single constraint
violation cascades through the network, with illustrative amplification
ratios on the order of \(10^3\) to \(10^4\) under modest network
parameters, because trust recovery is gradual while trust destruction is
instantaneous. Cooperation strictly dominates conflict for all
physically realizable friction parameters, providing the
energy-denominated cost structure that enters the game-theoretic payoff
matrices of the subsequent paper in this series.
\end{abstract}
\section{Introduction}\label{introduction}

\emph{Necessary Constraints} \citep{TC-I} proves that mutual boundary
constraints are necessary for the persistence of multiple agents in a
shared, finite-resource environment. Removing these constraints leads to
resource collision and conflict. This paper quantifies the thermodynamic
cost of that conflict.

When agents compete for a contested resource without cooperative
constraints, the interaction takes the form of a contest: each agent
expends energy to increase its probability of capturing the resource,
and the contest itself dissipates energy that neither agent recovers. We
model this using Tullock contest theory \citep{Tullock1980} with
physical parameters (boundary integrity, damage fractions, repair
multipliers) and prove that the total energy cost of conflict exceeds
the value of the contested resource under realistic conditions ---
destructive boundary damage between two agents, or any positive damage
among three or more.

The specific contributions of this paper are:

\begin{enumerate}
\def\labelenumi{\arabic{enumi}.}
\item
  \textbf{Contest equilibria} (Proposition~\ref{prop-symmetric-contest},
  Proposition~\ref{prop-asymmetric-contest}). We derive the symmetric
  and asymmetric Nash equilibria of the resource contest, showing that
  both agents expend energy proportional to the resource value, and that
  total contest dissipation approaches the full resource value as the
  number of contestants grows (Theorem~\ref{thm-dissipation-scaling}).
\item
  \textbf{Net-negative conflict}
  (Theorem~\ref{thm-net-negative-conflict}). We define the friction
  multiplier
  \(\Phi = 1 + \kappa(\psi^{\text{off}} + \psi^{\text{def}})\)
  incorporating contest expenditure, boundary damage, and
  thermodynamically irreversible repair costs, and prove that the total
  energy dissipated by conflict exceeds the value of the contested
  resource once boundary damage is sufficiently destructive
  (\(\kappa\psi > 1\)) or three or more agents contest it.
\item
  \textbf{Cooperation dominance}
  (Theorem~\ref{thm-cooperation-dominance}). We prove that the
  cooperative regime strictly dominates the conflict regime for all
  physically realizable friction parameters, providing the energy-cost
  foundation for the game-theoretic equilibrium results of
  \emph{Cooperative Equilibrium} \citep{TC-V}.
\item
  \textbf{Network cascading} (Theorem~\ref{thm-cascading-friction}). We
  prove that a single constraint violation in a cooperative network
  injects friction that degrades all future interactions, with
  illustrative amplification ratios of roughly \(10^3\) to \(10^4\)
  under modest network parameters, formalizing the asymmetry between
  trust destruction and trust recovery.
\end{enumerate}

All formal results are stated as numbered theorems with full proofs,
denominated in energy units, and computationally verified.

The paper proceeds as follows: model and contest framework, boundary
integrity, contest equilibria and full cost of conflict, the
net-negative theorem, cooperative versus conflict regimes, cascading
network effects, the interaction cost function, worked examples,
decisive contests and escalation dynamics, and discussion.

\subsection{Related Work}\label{related-work}

The mathematical machinery of this work (contest theory, Nash
equilibria, and boundary-damage models) draws on established literatures
in conflict economics, non-equilibrium thermodynamics, and network
theory. Our contribution is the synthesis: connecting contest-theoretic
dissipation to thermodynamic boundary costs and network friction to
derive a unified, energy-denominated cost function for adversarial
interaction.

\textbf{Contest theory and conflict economics.} \citet{Tullock1980}
introduced the ratio-form contest success function in the context of
rent-seeking, showing that competitions for rents dissipate a
substantial fraction of the rent's value. \citet{Skaperdas1996}
axiomatized the Tullock contest success function, proving it is the
unique function satisfying homogeneity, monotonicity, and independence
from irrelevant alternatives. \citet{Hirshleifer1991} analyzed conflict
as an economic activity, modeling the ``technology of conflict'' through
contest production functions and demonstrating that resources devoted to
fighting are pure waste from the system's perspective.
\citet{GarfinkelSkaperdas2012} provide a comprehensive treatment of the
economics of peace and conflict, synthesizing decades of work on arms
races, territorial disputes, and the welfare costs of insecurity.
\citet{Konrad2009} surveys the theory of contests more broadly, covering
multi-player, multi-stage, and asymmetric extensions. This literature
establishes that contest dissipation is a pervasive feature of
competitive interaction, but treats payoffs in dimensionless utility
units and does not connect dissipation to physical energy budgets. Our
framework denominates all contest payoffs in energy units and augments
the standard Tullock model with boundary-damage and repair costs derived
from thermodynamic principles.

\textbf{Non-equilibrium thermodynamics and boundary maintenance.} The
physical model of agents as open thermodynamic systems maintaining
low-entropy internal states draws on \citet{Schrodinger1944}, who
identified that living systems persist by importing free energy and
exporting entropy. \citet{Prigogine1978} demonstrated that systems far
from equilibrium can spontaneously generate ordered structures through
symmetry-breaking instabilities, but did not formalize the boundaries
that maintain organized entities as distinct agents or the energetic
cost of boundary breaches. \citet{England2013} derived a
statistical-mechanical lower bound on heat production during
self-replication. \citet{Friston2013} formalized persistence via Markov
blankets under the Free Energy Principle, addressing the single-agent
problem of what one entity does to maintain itself. \citet{Maturana1980}
described the operational closure of self-producing systems
(autopoiesis), the topology of self-maintenance that our boundary
integrity parameter \(B_i\) quantifies energetically. The thermodynamic
tradition establishes that boundary maintenance is energetically costly
and that entropy increase is thermodynamically irreversible; our
contribution is proving that boundary damage from conflict incurs repair
costs strictly exceeding the original damage (a consequence of the
Second Law), making conflict a thermodynamically amplified loss.

\textbf{Network effects and transaction costs.} The cascading friction
analysis draws on network theory and transaction cost economics.
\citet{Jackson2008} develops the theory of network formation and
stability, showing that network structure affects the efficiency of
economic exchange. \citet{Ostrom1990} documented how communities develop
institutional constraint systems (monitoring, graduated sanctions,
conflict resolution mechanisms) to manage shared resources, institutions
whose overhead costs correspond to our network friction coefficient
\(\phi\). The ecological economics tradition
\citep{Lotka1922, GeorgescuRoegen1971} recognized that energy is the
fundamental currency of biological and economic activity, and that the
Second Law imposes irreversible costs on all transformations
\citep{GeorgescuRoegen1979}. Our contribution is quantifying the
cascading cost of trust violations across cooperative networks and
proving that the network-level damage exceeds the direct bilateral cost
by orders of magnitude.

\section{Model}\label{model}

\subsection{Agents and Resources}\label{agents-and-resources}

We consider \(N \geq 2\) agents sharing an environment with \(n\)
resource dimensions and finite endowment
\(\mathbf{R} = (R_1, \ldots, R_n) \in \mathbb{R}_{>0}^n\). Each agent
\(i\) is an open thermodynamic system maintaining a boundary between its
internal low-entropy state and the external environment. Agent \(i\)
selects a strategy vector \(\mathbf{x}_i \in \mathbb{R}_{\geq 0}^n\),
and the physically realizable joint strategy must satisfy resource
conservation: \(\sum_{i=1}^N x_{ij} \leq R_j\) for all \(j\).

Each agent's utility function \(U_i(\mathbf{x}_i)\) satisfies the
regularity conditions of \emph{Necessary Constraints} \citep{TC-I}:
twice continuously differentiable, strictly concave, monotonically
increasing at the origin, and coercive. Under these conditions, each
agent has a unique unconstrained optimum \(\mathbf{x}_i^{\circ}\). When
\(N\) meets the collision threshold for at least one resource, the
aggregate unconstrained demand exceeds supply
(\(\sum_i x_{ij}^{\circ} > R_j\)), and no feasible joint allocation can
simultaneously satisfy all agents' unconstrained optima.

\emph{Necessary Constraints} \citep{TC-I} proves that boundary
constraints are necessary for feasible allocation in this setting. This
paper addresses the complementary question: what is the energetic cost
when those constraints are absent and the resulting collision is
resolved by direct resource competition?

\subsection{Boundary Integrity}\label{boundary-integrity}

Each agent \(i\) maintains a boundary characterized by its integrity
\(B_i > 0\), measured in energy units, representing the total energetic
investment in maintaining the separation between the agent's internal
low-entropy state and the external environment. Maintaining boundary
integrity against environmental entropy requires continuous work
\(C_{\text{maintain},i} = \gamma_i B_i\), where \(\gamma_i > 0\) is the
entropy leakage rate. This is the baseline metabolic cost of
persistence, independent of any external threat.

When an external agent breaches agent \(i\)'s boundary, the work
required scales with the boundary's integrity: full breach requires work
equal to \(B_i\). The specific resistance profile (uniform, hardening,
or shell) affects only the dynamics of partial breach, not the total
cost.

\subsection{Contest Framework}\label{contest-framework}

When two or more agents claim the same resource unit and no constraint
directs the allocation, the resource is resolved by a contest: each
agent invests energy to increase its probability of winning. We adopt
the Tullock contest success function \citep{Tullock1980, Skaperdas1996}:

\[p_A(e_A, e_B) = \frac{\sigma_A \, e_A^r}{\sigma_A \, e_A^r + \sigma_B \, e_B^r}\]

where \(e_i \geq 0\) is agent \(i\)'s contest investment (energy
diverted from productive use into conflict), \(\sigma_i > 0\) is
conflict effectiveness, and \(r > 0\) is the decisiveness parameter.
This is the unique contest success function satisfying homogeneity,
monotonicity, and independence from irrelevant alternatives
\citep{Skaperdas1996}.

Beyond the direct contest expenditure, conflict inflicts boundary damage
on both parties (offensive self-damage from diverting resources,
defensive damage from the opponent's attack), and repairing that damage
is thermodynamically more expensive than the original maintenance cost,
a direct consequence of the Second Law. The friction multiplier \(\Phi\)
captures the total cost amplification from these combined effects.

\section{The Boundary Integrity
Model}\label{the-boundary-integrity-model}

\subsection{Boundary as an Energetic
Structure}\label{boundary-as-an-energetic-structure}

Each agent \(i\) maintains a \textbf{boundary} (Markov blanket in
statistical physics; cell membrane in biology; legal/military defense in
sociology) characterized by its \textbf{integrity}:

\begin{definition}[Boundary Integrity]
\label{def-boundary-integrity}

Agent $i$'s boundary integrity is a scalar $B_i > 0$ measured in energy units, representing the total energetic investment in maintaining the boundary between the agent's internal low-entropy state and the external environment.
\end{definition}

\textbf{Physical interpretation:} \(B_i\) is the stored energy in the
agent's defensive infrastructure (the immune system, the cell wall, the
military budget, the information-security apparatus). A higher \(B_i\)
means a more robust boundary that is costlier to breach.

\subsection{Boundary Maintenance Cost}\label{boundary-maintenance-cost}

Maintaining boundary integrity against environmental entropy requires
continuous work:

\[C_{\text{maintain},i} = \gamma_i \cdot B_i\]

where \(\gamma_i > 0\) is the \textbf{entropy leakage rate}, the rate at
which the boundary degrades without active maintenance. This is the
baseline metabolic cost of persistence (maintaining identity),
independent of any external threat.

\subsection{Boundary Resistance
Function}\label{boundary-resistance-function}

When an external agent attempts to breach agent \(i\)'s boundary, the
boundary exerts a \textbf{resistance} that the attacker must overcome.
We model this as a force-displacement relationship:

\begin{definition}[Boundary Resistance]
\label{def-boundary-resistance}

The work required to breach agent $i$'s boundary to depth $d \in [0, 1]$ (where $d = 0$ is the exterior and $d = 1$ is full breach) is:

$$W_{\text{breach}}(d) = \int_0^d F_i(s) \, ds$$

where $F_i(s) = B_i \cdot f(s)$ is the resistance force and $f : [0, 1] \to \mathbb{R}_{> 0}$ is a normalized resistance profile with $\int_0^1 f(s) \, ds = 1$.
\end{definition}

The full breach cost is:

\[W_{\text{breach}}(1) = B_i \cdot \int_0^1 f(s) \, ds = B_i\]

That is, fully breaching agent \(i\)'s boundary requires work equal to
the boundary's integrity, an energetically intuitive result.

\textbf{Resistance profile choices:}

{\def\LTcaptype{none} %
\begin{longtable}[]{@{}
  >{\raggedright\arraybackslash}p{(\linewidth - 4\tabcolsep) * \real{0.3333}}
  >{\raggedright\arraybackslash}p{(\linewidth - 4\tabcolsep) * \real{0.3333}}
  >{\raggedright\arraybackslash}p{(\linewidth - 4\tabcolsep) * \real{0.3333}}@{}}
\toprule\noalign{}
\begin{minipage}[b]{\linewidth}\raggedright
Profile
\end{minipage} & \begin{minipage}[b]{\linewidth}\raggedright
\(f(s)\)
\end{minipage} & \begin{minipage}[b]{\linewidth}\raggedright
Physical meaning
\end{minipage} \\
\midrule\noalign{}
\endhead
\bottomrule\noalign{}
\endlastfoot
Uniform & \(f(s) = 1\) & Constant resistance (simple barrier) \\
Hardening & \(f(s) = 2s\) & Resistance increases with depth (layered
defense) \\
Shell & \(f(s) = 2(1 - s)\) & Hard exterior, soft interior (biological
cell) \\
\end{longtable}
}

The qualitative results below hold for any positive, integrable \(f\).
The specific profile affects only the dynamics of partial breach, not
the total cost.

\section{The Contest Model: Conflict as Resource
Competition}\label{the-contest-model-conflict-as-resource-competition}

\subsection{Setup}\label{setup}

Consider a single contested resource of value \(E_j > 0\) (energy units)
that is in collision (Definition 2 \citep{TC-I}). Two agents \(A\) and
\(B\) both claim this resource. In the absence of a boundary constraint
directing the allocation, the resource is resolved by a
\textbf{contest}: each agent invests energy to increase its probability
of winning.

\begin{definition}[Conflict Contest]
\label{def-conflict-contest}

A conflict contest over resource $E_j$ between agents $A$ and $B$ is characterized by:

(a) \textbf{Contest investments} $e_A, e_B \geq 0$: the energy each agent diverts from productive use into conflict.

(b) \textbf{Contest success function} $p_A(e_A, e_B) \in [0, 1]$: the probability that agent $A$ wins (secures) the resource, with $p_B = 1 - p_A$.

(c) \textbf{Payoffs:}
$$\Pi_A(e_A, e_B) = p_A(e_A, e_B) \cdot E_j - e_A$$
$$\Pi_B(e_A, e_B) = (1 - p_A(e_A, e_B)) \cdot E_j - e_B$$
\end{definition}

\subsection{The Tullock Contest Success
Function}\label{the-tullock-contest-success-function}

We adopt the axiomatically grounded \textbf{Tullock contest success
function} \citep{Skaperdas1996}:

\[p_A(e_A, e_B) = \frac{\sigma_A \, e_A^r}{\sigma_A \, e_A^r + \sigma_B \, e_B^r}\]

where:

\begin{itemize}
\tightlist
\item
  \(\sigma_i > 0\) is agent \(i\)'s \textbf{conflict effectiveness}
  (strength, weaponry, information advantage)
\item
  \(r > 0\) is the \textbf{decisiveness parameter}: \(r < 1\) means the
  contest is noisy (randomness dominates); \(r = 1\) is the baseline
  lottery; \(r > 1\) means the contest is decisive (larger investment
  almost certainly wins)
\end{itemize}

For \(e_A = e_B = 0\), we define
\(p_A = \sigma_A / (\sigma_A + \sigma_B)\) by convention.

\textbf{Axiomatization \citep{Skaperdas1996}:} This is the unique
contest success function satisfying: (i) Homogeneity of degree zero in
\((e_A, e_B)\) up to scaling; (ii) Monotonicity
(\(\partial p_A / \partial e_A > 0\),
\(\partial p_A / \partial e_B < 0\)); (iii) Independence from irrelevant
alternatives.

\subsection{Nash Equilibrium of the
Contest}\label{nash-equilibrium-of-the-contest}

Each agent chooses its contest investment to maximize its payoff, taking
the other's investment as given.

\textbf{Case 1: Symmetric Agents (\(\sigma_A = \sigma_B = 1\),
\(r = 1\))}

Agent \(A\)'s first-order condition:

\[\frac{\partial \Pi_A}{\partial e_A} = \frac{e_B}{(e_A + e_B)^2} \cdot E_j - 1 = 0\]

By symmetry, \(e_A^* = e_B^*\). Substituting:

\[\frac{e_A^*}{(2 e_A^*)^2} \cdot E_j = 1 \implies e_A^* = \frac{E_j}{4}\]

\begin{proposition}[Symmetric Contest Equilibrium]
\label{prop-symmetric-contest}

In a symmetric Tullock contest ($\sigma_A = \sigma_B$, $r = 1$) over resource $E_j$, the unique Nash equilibrium contest investments are:

$$e_A^* = e_B^* = \frac{E_j}{4}$$

Total contest dissipation:

$$D_2 = e_A^* + e_B^* = \frac{E_j}{2}$$

Each agent's expected payoff:

$$\Pi_A^* = \Pi_B^* = \frac{E_j}{2} - \frac{E_j}{4} = \frac{E_j}{4}$$
\end{proposition}

\textbf{Interpretation:} Half the resource value is consumed by the act
of fighting. Each agent expects to net only one quarter of the resource
value, compared to the cooperative outcome (splitting the resource, per
the Lagrangian constraints derivation) where each receives half with
zero contest expenditure.

\textbf{Case 2: Asymmetric Agents (\(V_A \neq V_B\), \(r = 1\))}

For the general case with potentially different valuations \(V_A\) and
\(V_B\) (how much each agent values the resource; we set
\(V_A = V_B = E_j\) for a common-value contest, but keep it general for
the derivation):

First-order conditions:

\[\frac{e_B}{(e_A + e_B)^2} \cdot V_A = 1, \qquad \frac{e_A}{(e_A + e_B)^2} \cdot V_B = 1\]

Dividing: \(\frac{e_B}{e_A} = \frac{V_B}{V_A}\), so
\(e_B = \frac{V_B}{V_A} e_A\).

Substituting back into \(A\)'s FOC:

\[\frac{V_B e_A / V_A}{e_A^2 (1 + V_B/V_A)^2} \cdot V_A = 1\]

\[e_A^* = \frac{V_A^2 V_B}{(V_A + V_B)^2}, \qquad e_B^* = \frac{V_A V_B^2}{(V_A + V_B)^2}\]

\begin{proposition}[Asymmetric Contest Equilibrium]
\label{prop-asymmetric-contest}

In a Tullock contest ($r = 1$) with valuations $V_A, V_B > 0$:

$$e_A^* = \frac{V_A^2 V_B}{(V_A + V_B)^2}, \qquad e_B^* = \frac{V_A V_B^2}{(V_A + V_B)^2}$$

Total dissipation:

$$D_2 = e_A^* + e_B^* = \frac{V_A V_B}{V_A + V_B}$$

This is half the harmonic mean of the valuations. For $V_A = V_B = E_j$: $D_2 = E_j/2$.
\end{proposition}

\textbf{Note:} The total dissipation \(D_2 = V_A V_B / (V_A + V_B)\) is
increasing in each valuation separately
(\(\partial D_2/\partial V_A = V_B^2/(V_A+V_B)^2 > 0\)), so it has no
unconstrained interior maximum. For a \emph{fixed aggregate valuation}
\(V_A + V_B\), however, \(D_2\) is maximized when \(V_A = V_B\) (equal
opponents waste the most) and vanishes as \(V_B \to 0\). When one agent
values the resource much more (\(V_B \ll V_A\), holding the aggregate
fixed), the weaker agent invests very little and the contest is
``cheap'', but the outcome is effectively determined by power asymmetry,
which corresponds to exploitative resource extraction.

\subsection{N-Agent Contest: Scaling to
Populations}\label{n-agent-contest-scaling-to-populations}

When \(N \geq 2\) symmetric agents (\(\sigma_i = 1\), \(V_i = E_j\),
\(r = 1\)) contest the same resource:

By symmetry, each invests \(e^*\) and wins with probability \(1/N\).
Agent \(i\)'s FOC:

\[\frac{(N-1) e^*}{N^2 (e^*)^2} \cdot E_j = 1 \implies e^* = \frac{(N-1)}{N^2} \cdot E_j\]

\begin{theorem}[Dissipation Scaling]
\label{thm-dissipation-scaling}

In a symmetric $N$-agent Tullock contest ($r = 1$) over resource $E_j$:

(a) Each agent's equilibrium investment:

$$e_i^* = \frac{N - 1}{N^2} \cdot E_j$$

(b) Total contest dissipation:

$$D_N = N \cdot e_i^* = \frac{N - 1}{N} \cdot E_j$$

(c) Each agent's expected payoff:

$$\Pi_i^* = \frac{E_j}{N} - \frac{(N-1) E_j}{N^2} = \frac{E_j}{N^2}$$

(d) Dissipation ratio:

$$\frac{D_N}{E_j} = \frac{N - 1}{N} = 1 - \frac{1}{N}$$
\end{theorem}

\begin{proof}
Part (a): With $N$ symmetric agents, agent $i$'s success probability is:

$$p_i = \frac{e_i}{e_i + \sum_{k \neq i} e_k}$$

At a symmetric equilibrium ($e_k = e^*$ for all $k$), $p_i = 1/N$. The FOC is:

$$\frac{\partial \Pi_i}{\partial e_i} = \frac{(N-1) e^*}{(N e^*)^2} \cdot E_j - 1 = 0$$

Solving: $e^* = (N-1) E_j / N^2$. Parts (b)–(d) follow by direct computation.
\end{proof}

\begin{corollary}[Vanishing Individual Returns]
\label{cor-vanishing-returns}

$$\lim_{N \to \infty} \Pi_i^* = \lim_{N \to \infty} \frac{E_j}{N^2} = 0$$

As the number of contestants grows, each agent's expected return collapses to zero while the aggregate waste approaches the full resource value. The contest becomes a pure value-destruction mechanism.
\end{corollary}

\begin{corollary}[Unit Dissipation Limit]
\label{cor-unit-dissipation-limit}

$$\lim_{N \to \infty} \frac{D_N}{E_j} = 1$$

In the limit of many contestants, the \textbf{entire resource value is dissipated in conflict}. No value remains for any agent. This is the mathematical expression of the "tragedy of the commons" reframed as a thermodynamic result.
\end{corollary}

{\def\LTcaptype{none} %
\begin{longtable}[]{@{}
  >{\raggedright\arraybackslash}p{(\linewidth - 8\tabcolsep) * \real{0.2000}}
  >{\raggedright\arraybackslash}p{(\linewidth - 8\tabcolsep) * \real{0.2000}}
  >{\raggedright\arraybackslash}p{(\linewidth - 8\tabcolsep) * \real{0.2000}}
  >{\raggedright\arraybackslash}p{(\linewidth - 8\tabcolsep) * \real{0.2000}}
  >{\raggedright\arraybackslash}p{(\linewidth - 8\tabcolsep) * \real{0.2000}}@{}}
\toprule\noalign{}
\begin{minipage}[b]{\linewidth}\raggedright
\(N\)
\end{minipage} & \begin{minipage}[b]{\linewidth}\raggedright
Individual investment \(e_i^*\)
\end{minipage} & \begin{minipage}[b]{\linewidth}\raggedright
Total dissipation \(D_N\)
\end{minipage} & \begin{minipage}[b]{\linewidth}\raggedright
\(D_N / E_j\)
\end{minipage} & \begin{minipage}[b]{\linewidth}\raggedright
Individual payoff \(\Pi_i^*\)
\end{minipage} \\
\midrule\noalign{}
\endhead
\bottomrule\noalign{}
\endlastfoot
2 & \(E_j/4\) & \(E_j/2\) & 50\% & \(E_j/4\) \\
3 & \(2E_j/9\) & \(2E_j/3\) & 66.7\% & \(E_j/9\) \\
5 & \(4E_j/25\) & \(4E_j/5\) & 80\% & \(E_j/25\) \\
10 & \(9E_j/100\) & \(9E_j/10\) & 90\% & \(E_j/100\) \\
100 & \(99E_j/10000\) & \(99E_j/100\) & 99\% & \(E_j/10000\) \\
\end{longtable}
}

\textbf{Physical reading:} This table demonstrates that the energy spent
fighting exceeds the energy gained from the resource. Even with just 2
agents, half the value is burned. With 10 agents, 90\% is burned. The
resource does not get ``won''; it gets \textbf{lost}.

\section{The Full Cost of Conflict: Boundary Damage and
Repair}\label{the-full-cost-of-conflict-boundary-damage-and-repair}

\subsection{Beyond Contest
Expenditure}\label{beyond-contest-expenditure}

The Tullock contest model captures the energy directly invested in
fighting (the ``ammunition''), but conflict inflicts additional costs:

\begin{enumerate}
\def\labelenumi{\arabic{enumi}.}
\tightlist
\item
  \textbf{Offensive boundary damage:} the attacker's boundary degrades
  from neglect (energy diverted to attack) and counter-damage.
\item
  \textbf{Defensive boundary damage:} the defender's boundary is
  physically damaged by the attack.
\item
  \textbf{Repair costs:} both agents must subsequently invest energy to
  restore boundary integrity.
\end{enumerate}

These costs are \textbf{on top of} the contest dissipation \(D_N\)
computed above.

\subsection{Boundary Damage Functions}\label{boundary-damage-functions}

\begin{definition}[Offensive Damage]
\label{def-offensive-damage}

In a conflict where agent $i$ invests $e_i$ in the contest, the \textbf{offensive damage} to the attacker's own boundary is:

$$\Delta B_i^{\text{off}} = \psi^{\text{off}} \cdot e_i$$

where $\psi^{\text{off}} \in [0, 1]$ is the offensive damage coefficient, the fraction of contest investment that translates into self-boundary degradation (energy diverted from maintenance, physical exposure during attack, counter-strikes).
\end{definition}

\begin{definition}[Defensive Damage]
\label{def-defensive-damage}

In a conflict where agent $i$ invests $e_i$ in the contest, the \textbf{defensive damage} inflicted on the opponent's boundary is:

$$\Delta B_j^{\text{def}} = \psi^{\text{def}} \cdot e_i$$

where $\psi^{\text{def}} \in [0, 1]$ is the defensive damage coefficient, the fraction of agent $i$'s attack energy that penetrates and damages agent $j$'s boundary.
\end{definition}

\textbf{Interpretation:} In a two-agent conflict where both invest
\(e^*\):

\begin{itemize}
\tightlist
\item
  Agent \(A\)'s total boundary damage:
  \(\Delta B_A = \psi^{\text{off}} \cdot e_A^* + \psi^{\text{def}} \cdot e_B^*\)
\item
  Agent \(B\)'s total boundary damage:
  \(\Delta B_B = \psi^{\text{off}} \cdot e_B^* + \psi^{\text{def}} \cdot e_A^*\)
\end{itemize}

For symmetric agents (\(e_A^* = e_B^* = e^*\)):

\[\Delta B_A = \Delta B_B = (\psi^{\text{off}} + \psi^{\text{def}}) \cdot e^*\]

\subsection{Repair Costs: The Thermodynamic
Asymmetry}\label{repair-costs-the-thermodynamic-asymmetry}

Repairing a damaged boundary is thermodynamically more expensive than
maintaining an intact one. This is a direct consequence of the Second
Law: destruction increases entropy; reconstruction requires importing
negentropy (ordered energy) from the environment. To make the repair
multiplier a grounded ratio rather than a stipulated one, we first fix
\emph{what energy} the symbols \(B_i\) and \(\Delta B_i\) denote.

\begin{definition}[Embodied Boundary Energy]
\label{def-embodied-energy}

The boundary integrity $B_i$ (Definition~\ref{def-boundary-integrity}) is the \textbf{embodied (negentropic) energy} of the boundary: the total ordered free energy that must be supplied to assemble it from dispersed environmental matter into its functional low-entropy configuration,

$$B_i = \Delta G_{\text{boundary}} = \Delta H - T_0 \, \Delta S,$$

where $\Delta H$ is the enthalpic (bond/configurational) content, $\Delta S < 0$ is the entropy reduction of ordering, and $T_0$ is the environmental temperature. The \textbf{boundary damage} $\Delta B_i$ is the embodied energy of the destroyed portion of the boundary, measured in the same energy units as $B_i$: the ordered free energy that has been dispersed and must be re-supplied to restore the pre-damage configuration.
\end{definition}

This is not an additional axiom. It is the physical content of the
normalization \(W_{\text{breach}}(1) = B_i\) established above: fully
breaching a boundary dismantles its ordered structure, and the minimum
work to dismantle equals the free energy stored in it, so the resistance
profile \(f(s)\) only distributes \(B_i\) along the depth coordinate
(\(\int_0^1 f = 1\) guarantees the total is \(B_i\)). The maintenance
flow \(C_{\text{maintain},i} = \gamma_i B_i\) is the depreciation on
this stock; \(B_i\) is the capital.

\begin{definition}[Repair Multiplier]
\label{def-repair-multiplier}

The energy required to repair boundary damage $\Delta B_i$ is:

$$W_{\text{repair},i} = \kappa \cdot \Delta B_i$$

where $\kappa \geq 1$ is the \textbf{repair multiplier}. Restoring the destroyed portion means re-supplying its embodied energy $\Delta B_i$ \textbf{plus} the surcharges of doing so on a damaged, disordered site, so $\kappa$ decomposes as a surcharge over re-embodiment:

$$\kappa = 1 + \underbrace{\frac{W_{\text{clear}}}{\Delta B}}_{\kappa_{\text{clear}}} + \underbrace{\frac{W_{\text{penalty}}}{\Delta B}}_{\kappa_{\text{pen}}} + \underbrace{\frac{T_0 \, \Delta S_{\text{irr}}}{\Delta B}}_{\kappa_{\text{irr}}}.$$

The leading $1$ is the irreducible re-embodiment: at minimum one must re-supply what was lost. Every surcharge is non-negative, giving $\kappa \geq 1$; the strict inequality $\kappa > 1$ holds in any physical system because:

(i) Repair requires first clearing the damage (entropy removal, the term $\kappa_{\text{clear}}$), then rebuilding (negentropy import), two thermodynamic steps vs. one for maintenance.

(ii) The damaged state has higher entropy than the pre-damage state; returning to the original configuration requires at least $\Delta S \cdot T$ additional work (the term $\kappa_{\text{irr}}$; Clausius inequality / minimum work theorem, Second Law of Thermodynamics).

The term $\kappa_{\text{pen}}$ captures rebuilding under conditions worse than the original build (degraded site, lost economies of scale, matching to surviving structure).
\end{definition}

\textbf{The effort-linked damage form as a transfer relation.} The
damage functions \(\Delta B_i^{\text{off}} = \psi^{\text{off}} e_i\) and
\(\Delta B_j^{\text{def}} = \psi^{\text{def}} e_i\) introduced below tie
the embodied loss to the attacker's contest effort through a
\emph{transfer coefficient}
\(\psi = \chi \cdot \rho_B / \varepsilon_{\text{break}}\), the product
of a coupling fraction \(\chi \in [0,1]\) (how much effort reaches the
boundary) and a leverage ratio \(\rho_B / \varepsilon_{\text{break}}\)
(embodied energy freed per unit fracture energy delivered). For any
structure cheaper to break than to build, the leverage ratio exceeds
unity: this is where the ``cheap to destroy'' asymmetry lives. It is
booked in \(\psi\), not in \(\kappa\) --- \(\kappa\) measures repair
cost relative to the embodied energy lost, a separate quantity that
stays \(O(1)\) even when destruction is highly leveraged. The bounded
coefficients \(\psi^{\text{off}}, \psi^{\text{def}} \in [0,1]\) used
below are the coupling fraction \(\chi\) under the normalization
\(\rho_B \approx \varepsilon_{\text{break}}\).

\textbf{Empirical calibration.} The quantity \(\kappa\) measures is an
\emph{energy ratio}: the metabolic energy expended to regenerate damaged
structure divided by the embodied energy of the structure restored ---
not a metabolic \emph{rate} multiplier. In biological healing,
bioenergetic accounting places net growth efficiency at roughly
60--75\%, so depositing one unit of tissue energy costs about 1.2--1.7
units of metabolic energy \citep{Reeds1982EnergyCosts}, a synthesis
surcharge over re-embodiment (the
\(\kappa_{\text{clear}} + \kappa_{\text{pen}}\) terms). This 1.2--1.7
interval is the range that re-embodiment energetics \emph{cleanly}
support as an energy ratio. Wound healing adds an inflammatory surge on
top; clinically this appears as an elevated basal metabolic \emph{rate}
--- a power draw raised by factors of roughly 1.4 to 2.3 during major
healing \citep{Demling2009} --- which is not the ratio \(\kappa\)
itself, but which, integrated over the healing episode, drives the
\emph{total} energy overhead above the clean synthesis range. We
therefore take \(\kappa = 2\) as a deliberately conservative upper
baseline in worked examples: it sits \emph{above} the clean
synthesis-ratio range of 1.2--1.7, the excess folding in the integrated
inflammatory overhead of a full healing episode rather than resting on
the synthesis ratio alone. We report the sensitivity of the
net-negativity verdict around the sign-flip line \(\kappa\psi = 1\)
below.

\subsection{Total Conflict Cost: The Complete
Ledger}\label{total-conflict-cost-the-complete-ledger}

We now assemble all cost components for a two-agent symmetric contest
(\(r = 1\), \(\sigma_A = \sigma_B\)). Writing
\(\psi \triangleq \psi^{\text{off}} + \psi^{\text{def}}\) for the
\textbf{combined boundary-damage fraction}, each agent's per-episode
cost is its contest investment plus the repair that restores its
boundary:

{\def\LTcaptype{none} %
\begin{longtable}[]{@{}
  >{\raggedright\arraybackslash}p{(\linewidth - 4\tabcolsep) * \real{0.3333}}
  >{\raggedright\arraybackslash}p{(\linewidth - 4\tabcolsep) * \real{0.3333}}
  >{\raggedright\arraybackslash}p{(\linewidth - 4\tabcolsep) * \real{0.3333}}@{}}
\toprule\noalign{}
\begin{minipage}[b]{\linewidth}\raggedright
Cost component
\end{minipage} & \begin{minipage}[b]{\linewidth}\raggedright
Per agent
\end{minipage} & \begin{minipage}[b]{\linewidth}\raggedright
System total
\end{minipage} \\
\midrule\noalign{}
\endhead
\bottomrule\noalign{}
\endlastfoot
Contest investment \(e_i^*\) & \(E_j / 4\) & \(E_j / 2\) \\
Boundary repair \(W_{\text{repair},i} = \kappa \, \Delta B_i\) &
\(\kappa \, \psi \cdot E_j / 4\) & \(\kappa \, \psi \cdot E_j / 2\) \\
\textbf{Total cost per agent} & \([1 + \kappa \psi] \cdot E_j / 4\) & \\
\textbf{Total system cost} & & \([1 + \kappa \psi] \cdot E_j/2\) \\
\end{longtable}
}

The destroyed embodied energy \(\Delta B_i = \psi \, e_i^*\) does
\textbf{not} appear as a separate additive cost: the repair term
\(W_{\text{repair},i} = \kappa \Delta B_i\) already re-supplies it (the
leading \(1\) in \(\kappa\) is precisely the re-embodiment of
\(\Delta B_i\)), so charging \(\Delta B_i\) again would double-book the
destroyed structure. Counting a standalone \(\Delta B_i\) \emph{and} the
full \(\kappa \Delta B_i\) would inflate the per-agent multiplier from
the correct \(1+\kappa\psi\) to \(1 + (1+\kappa)\psi\) --- an over-count
of exactly the destroyed \(\Delta B_i = \psi\, e_i^*\) per agent; the
correct ledger is contest effort plus repair.

Define the \textbf{total friction multiplier}:

\[\Phi \triangleq 1 + \kappa \, \psi = 1 + \kappa(\psi^{\text{off}} + \psi^{\text{def}})\]

Then:

\[L_{\text{system}}^{(2)} = \Phi \cdot \frac{E_j}{2}\]

where \(L_{\text{system}}^{(2)}\) is the total system loss from a
two-agent contest.

\section{The Net-Negative Theorem}\label{the-net-negative-theorem}

\subsection{Statement}\label{statement}

\begin{theorem}[Net-Negative Conflict]
\label{thm-net-negative-conflict}

In a conflict contest over resource $E_j$ between $N \geq 2$ symmetric agents with friction multiplier $\Phi = 1 + \kappa(\psi^{\text{off}} + \psi^{\text{def}})$:

(a) The total system loss (energy dissipated by conflict) is:

$$L_{\text{system}}^{(N)} = \Phi \cdot \frac{N - 1}{N} \cdot E_j$$

(b) The net system value (resource obtained minus all costs) is:

$$V_{\text{net}} = E_j - L_{\text{system}}^{(N)} = E_j \left(1 - \Phi \cdot \frac{N - 1}{N}\right)$$

(c) The conflict is \textbf{net-negative} ($V_{\text{net}} < 0$) if and only if:

$$\Phi > \frac{N}{N - 1}$$

For $N = 2$: $\Phi > 2$ (equivalently $\kappa\psi > 1$). For $N = 3$: $\Phi > 3/2$; since $N/(N-1) \leq 3/2$ for all $N \geq 3$, the condition $\Phi > 3/2$ suffices for every $N \geq 3$. As $N \to \infty$ the threshold falls to $\Phi > 1$.
\end{theorem}

\begin{proof}
The contest dissipation for $N$ symmetric agents is $D_N = \frac{N-1}{N} E_j$ (Theorem~\ref{thm-dissipation-scaling}). Each agent's boundary damage is $\Delta B_i = \psi \, e_i^*$ with $\psi = \psi^{\text{off}} + \psi^{\text{def}}$, and restoring it costs $W_{\text{repair},i} = \kappa \Delta B_i = \kappa \psi \, e_i^*$ (which re-supplies the destroyed $\Delta B_i$, so damage is not counted separately). Summed over $N$ agents, using $N e_i^* = D_N$, the total repair cost is $\kappa \psi \cdot \frac{N-1}{N} E_j$. Adding contest dissipation:

$$L_{\text{system}}^{(N)} = \left[1 + \kappa(\psi^{\text{off}} + \psi^{\text{def}})\right] \cdot \frac{N-1}{N} E_j = \Phi \cdot \frac{N-1}{N} E_j$$

The net value is $V_{\text{net}} = E_j - L_{\text{system}}^{(N)}$. Setting $V_{\text{net}} < 0$:

$$E_j < \Phi \cdot \frac{N-1}{N} E_j \iff 1 < \Phi \cdot \frac{N-1}{N} \iff \Phi > \frac{N}{N-1}$$

All damage terms are aggregate: $L_{\text{system}}^{(N)}$ sums contest dissipation and repair across all $N$ agents, not per-pair. For asymmetric agents, the per-agent repair costs $\kappa(\psi_i^{\text{off}} + \psi_i^{\text{def}}) e_i^*$ are summed over all $i$.
\end{proof}

\subsection{Physical
Interpretation}\label{physical-interpretation-net-negative}

\begin{corollary}[Inevitability of Net-Negative Conflict]
\label{cor-inevitability-net-negative}

For any physical system where boundary damage is non-negligible ($\psi^{\text{off}} + \psi^{\text{def}} > 0$) and repair is thermodynamically irreversible ($\kappa > 1$), the conflict is net-negative for sufficiently many contestants.
\end{corollary}

\begin{proof}
As $N \to \infty$, the threshold $N/(N-1) \to 1$. Since $\Phi = 1 + \kappa(\psi^{\text{off}} + \psi^{\text{def}}) > 1$ whenever $\psi^{\text{off}} + \psi^{\text{def}} > 0$, there exists a finite $N^*$ such that $\Phi > N^*/(N^* - 1)$, beyond which the conflict is net-negative. Specifically:

$$N^* = \left\lfloor \frac{\Phi}{\Phi - 1} \right\rfloor + 1$$
\end{proof}

\begin{corollary}[The Two-Agent Threshold]
\label{cor-two-agent-threshold}

For two agents ($N = 2$), the conflict is net-negative when $\Phi > 2$, i.e., when:

$$\kappa \, \psi = \kappa(\psi^{\text{off}} + \psi^{\text{def}}) > 1$$

At the baseline $\kappa = 2$ this requires $\psi > 1/2$: the two-agent contest tips net-negative only once boundary damage is \emph{destructive}. At the mild baseline $\psi^{\text{off}} = 0.15$, $\psi^{\text{def}} = 0.25$ ($\psi = 0.4$), $\Phi = 1 + 2 \cdot 0.4 = 1.8 < 2$ and the two-agent contest is net-positive; at a destructive $\psi^{\text{off}} = \psi^{\text{def}} = 0.3$ ($\psi = 0.6$), $\Phi = 1 + 2 \cdot 0.6 = 2.2 > 2$ and it is net-negative. For three or more agents the threshold falls to $\Phi > N/(N-1) \leq 3/2$, so any positive boundary damage makes the contest net-negative (Theorem~\ref{thm-net-negative-conflict}).
\end{corollary}

\subsection{Numerical Illustration: The Friction Multiplier
Landscape}\label{numerical-illustration-the-friction-multiplier-landscape}

The following table shows the total system loss
\(L_{\text{system}}^{(N)} / E_j\) for various parameter combinations,
with \(\psi = \psi^{\text{off}} + \psi^{\text{def}}\) and
\(\Phi = 1 + \kappa\psi\):

{\def\LTcaptype{none} %
\begin{longtable}[]{@{}
  >{\raggedright\arraybackslash}p{(\linewidth - 12\tabcolsep) * \real{0.1429}}
  >{\raggedright\arraybackslash}p{(\linewidth - 12\tabcolsep) * \real{0.1429}}
  >{\raggedright\arraybackslash}p{(\linewidth - 12\tabcolsep) * \real{0.1429}}
  >{\raggedright\arraybackslash}p{(\linewidth - 12\tabcolsep) * \real{0.1429}}
  >{\raggedright\arraybackslash}p{(\linewidth - 12\tabcolsep) * \real{0.1429}}
  >{\raggedright\arraybackslash}p{(\linewidth - 12\tabcolsep) * \real{0.1429}}
  >{\raggedright\arraybackslash}p{(\linewidth - 12\tabcolsep) * \real{0.1429}}@{}}
\toprule\noalign{}
\begin{minipage}[b]{\linewidth}\raggedright
\(\psi = \psi^{\text{off}} + \psi^{\text{def}}\)
\end{minipage} & \begin{minipage}[b]{\linewidth}\raggedright
\(\kappa\)
\end{minipage} & \begin{minipage}[b]{\linewidth}\raggedright
\(\Phi\)
\end{minipage} & \begin{minipage}[b]{\linewidth}\raggedright
\(N = 2\)
\end{minipage} & \begin{minipage}[b]{\linewidth}\raggedright
\(N = 5\)
\end{minipage} & \begin{minipage}[b]{\linewidth}\raggedright
\(N = 10\)
\end{minipage} & \begin{minipage}[b]{\linewidth}\raggedright
\(N = 100\)
\end{minipage} \\
\midrule\noalign{}
\endhead
\bottomrule\noalign{}
\endlastfoot
0 (pure contest) & n/a & 1.0 & 50\% & 80\% & 90\% & 99\% \\
0.2 & 1.5 & 1.3 & 65\% & 104\% & 117\% & 128.7\% \\
0.4 & 2.0 & 1.8 & 90\% & 144\% & 162\% & 178.2\% \\
0.6 & 2.0 & 2.2 & 110\% & 176\% & 198\% & 217.8\% \\
0.8 & 3.0 & 3.4 & 170\% & 272\% & 306\% & 336.6\% \\
\end{longtable}
}

Values exceeding 100\% represent \textbf{net-negative outcomes}, where
the system dissipates more energy than the contested resource contains.
With modest boundary damage parameters (\(\psi = 0.2\),
\(\kappa = 1.5\)), two-agent contests remain net-positive (dissipating
65\% of the resource value) while five-agent contests are already
net-negative. At the mild baseline (\(\psi = 0.4\), \(\kappa = 2\)) the
two-agent contest dissipates 90\% of the resource --- costly but still
net-positive --- while every contest with three or more agents is
net-negative (\(\Phi = 1.8 > 3/2\)).

Fixing the baseline \(\kappa = 2\) and reading across \(\psi\) isolates
the two-agent sign flip:

{\def\LTcaptype{none} %
\begin{longtable}[]{@{}
  >{\raggedright\arraybackslash}p{(\linewidth - 8\tabcolsep) * \real{0.2000}}
  >{\raggedright\arraybackslash}p{(\linewidth - 8\tabcolsep) * \real{0.2000}}
  >{\raggedright\arraybackslash}p{(\linewidth - 8\tabcolsep) * \real{0.2000}}
  >{\raggedright\arraybackslash}p{(\linewidth - 8\tabcolsep) * \real{0.2000}}
  >{\raggedright\arraybackslash}p{(\linewidth - 8\tabcolsep) * \real{0.2000}}@{}}
\toprule\noalign{}
\begin{minipage}[b]{\linewidth}\raggedright
Regime
\end{minipage} & \begin{minipage}[b]{\linewidth}\raggedright
\(\kappa\)
\end{minipage} & \begin{minipage}[b]{\linewidth}\raggedright
\(\psi\)
\end{minipage} & \begin{minipage}[b]{\linewidth}\raggedright
\(\Phi = 1 + \kappa\psi\)
\end{minipage} & \begin{minipage}[b]{\linewidth}\raggedright
Two-agent outcome (\(E_j = 100\))
\end{minipage} \\
\midrule\noalign{}
\endhead
\bottomrule\noalign{}
\endlastfoot
Mild dispute & 2 & 0.4 & 1.8 & Net-positive (\(+10\)) \\
Break-even & 2 & 0.5 & 2.0 & Zero (\(0\)) \\
Destructive & 2 & 0.6 & 2.2 & Net-negative (\(-10\)) \\
Any \(\psi > 0\), \(N \geq 3\) & 2 & 0.4 & 1.8 & Net-negative
(\(\Phi > 3/2\)) \\
\end{longtable}
}

\textbf{Sensitivity around the sign-flip line.} The two-agent verdict
turns on the single product \(\kappa\psi\) relative to \(1\), so it is
sensitive near \(\kappa\psi = 1\): at the baseline \(\kappa = 2\) a
swing in \(\psi\) from \(0.4\) to \(0.6\) moves the two-agent net from
\(+10\) to \(-10\) on \(E_j = 100\). The headline conclusion does not
rest on this knife-edge. Net-negativity is reached \textbf{either} by
destructive damage (\(\kappa\psi > 1\), e.g.~\(\psi \gtrsim 0.5\) at
\(\kappa = 2\)) \textbf{or} by three or more contestants (any positive
damage suffices once \(\Phi > N/(N-1)\)). The \(N\)-agent result and the
empirically observed conflict cost-to-value ratios --- which range from
roughly \(2.5{:}1\) to \(500{:}1\) across litigation, military, and
cyber conflict (Applications, below) --- carry the net-negativity claim
independently of the precise value of \(\kappa\), since those ratios
measure the full realized loss \(\Phi \cdot (N-1)/N\) rather than
\(\kappa\) alone.

\section{Comparison: Conflict vs.~Cooperative
Regimes}\label{comparison-conflict-vs.-cooperative-regimes}

\subsection{The Cooperative Baseline}\label{the-cooperative-baseline}

Under the boundary-constraint-based cooperative regime (Theorem 7
\citep{TC-I}), agents achieve the variational equilibrium allocation
with:

\begin{itemize}
\tightlist
\item
  \textbf{Zero contest expenditure:} No agent invests energy in
  fighting.
\item
  \textbf{Zero boundary damage:} No boundaries are breached.
\item
  \textbf{Zero repair costs:} No damage to repair.
\item
  \textbf{Shadow price costs only:} Each agent's cost is bounded by the
  Lagrange multipliers \(\mu_{Bk}^*\), the opportunity cost of
  respecting others' boundary constraints.
\end{itemize}

The total system welfare under cooperation (Theorem 7 \citep{TC-I}) is:

\[SW_{\text{coop}} = \sum_{i=1}^N U_i(\mathbf{x}_i^*)\]

\subsection{The Conflict Regime}\label{the-conflict-regime}

Under conflict (no boundary constraints), the total system welfare is:

\[SW_{\text{conflict}} = \sum_{i=1}^N \left[ p_i \cdot E_j - e_i^* - W_{\text{repair},i} \right] = E_j - D_N - W_{\text{repair, total}}\]

\[= E_j - \Phi \cdot \frac{N-1}{N} E_j = E_j \left(1 - \Phi \cdot \frac{N-1}{N}\right)\]

\subsection{The Cooperation Premium}\label{the-cooperation-premium}

\begin{theorem}[Cooperation Dominance]
\label{thm-cooperation-dominance}

Suppose utilities are energy-denominated, so that the cooperative allocation realizes at least the resource's energy value ($SW_{\text{coop}} \geq E_j$). Then the cooperative regime strictly dominates the conflict regime whenever resource collision occurs. Specifically:

$$SW_{\text{coop}} - SW_{\text{conflict}} \geq \Phi \cdot \frac{N-1}{N} \cdot E_j > 0$$

The cooperation premium is:
\begin{itemize}
\item Increasing in $N$ (more agents → more waste in conflict)
\item Increasing in $\Phi$ (more boundary damage → more destructive conflict)
\item Increasing in $E_j$ (higher-value resources → higher absolute waste)
\end{itemize}

\end{theorem}

\begin{proof}
Under cooperation, the full resource value $E_j$ is available for productive allocation (minus only the shadow-price opportunity costs, which are transfers between agents, not system losses); with energy-denominated utilities this yields $SW_{\text{coop}} \geq E_j$. Under conflict, $SW_{\text{conflict}} = E_j - L_{\text{system}}^{(N)}$ with $L_{\text{system}}^{(N)} = \Phi \cdot \frac{N-1}{N} E_j$ lost to pure waste. Therefore:

$$SW_{\text{coop}} - SW_{\text{conflict}} \geq E_j - \left(E_j - L_{\text{system}}^{(N)}\right) = L_{\text{system}}^{(N)} = \Phi \cdot \frac{N-1}{N} \cdot E_j > 0$$

The inequality is strict because $\Phi \geq 1$, $N \geq 2$, and $E_j > 0$. Absent the energy-denomination premise the same argument gives the weaker bound $SW_{\text{coop}} - SW_{\text{conflict}} \geq SW_{\text{coop}} - E_j + L_{\text{system}}^{(N)}$, which remains strictly positive whenever $SW_{\text{coop}} > E_j - L_{\text{system}}^{(N)}$.
\end{proof}

\textbf{Physical interpretation:} This is the formal proof of our
central claim: \emph{``The only way to globally minimize
\(E_{\text{conflict}}\) across the system is to establish mutual
constraints.''} The cooperative premium quantifies exactly how much
energy the system saves by adopting boundary constraints.

\subsection{Worked Example: Cooperation vs.~Conflict Under
Scarcity}\label{worked-example-cooperation-vs.-conflict-under-scarcity}

Consider two symmetric agents with utility parameters \(\alpha = 10\),
\(\beta = 1\) (matching the worked example of \emph{Necessary
Constraints} \citep{TC-I}), shared resource \(R = 12\), and friction
multiplier \(\Phi = 1.8\)
(\(\psi^{\text{off}} = \psi^{\text{def}} = 0.2\), \(\psi = 0.4\),
\(\kappa = 2\)):

\textbf{Cooperative regime (VE):} - Allocation: \(x_A^* = x_B^* = 6\) -
Utilities: \(U_A = U_B = 42\) - System welfare:
\(SW_{\text{coop}} = 84\)

\textbf{Conflict regime:}

The contested portion is the collision overlap. Each agent wants 10 but
only 12 is available. The contested amount is
\(\sum x_i^{\circ} - R = 20 - 12 = 8\) units.

Contest over the excess demand (\(E_{\text{contested}} = 8\) in resource
units, with energy value depending on the utility function):

\begin{itemize}
\tightlist
\item
  Contest dissipation: \(D_2 = E_{\text{contested}}/2 = 4\) (energy
  units lost to fighting)
\item
  Boundary repair: \(\kappa \psi \cdot D_2 = 2 \cdot 0.4 \cdot 4 = 3.2\)
  (the destroyed embodied energy is re-supplied within this term, not
  charged separately)
\item
  Total system loss: \(D_2 + 3.2 = 7.2\) energy units (equivalently
  \(\Phi \cdot D_2 = 1.8 \cdot 4\))
\end{itemize}

Even a simplified conflict model shows the cooperation premium: the
system saves at least 7.2 energy units by establishing mutual
constraints, nearly the whole of the contested portion of the resource.

\section{Cascading Network Effects: Trust and Ambient
Friction}\label{cascading-network-effects-trust-and-ambient-friction}

\subsection{Motivation}\label{motivation}

The preceding analysis treats each conflict as isolated. In reality,
conflict events propagate through the network: when agent \(A\) violates
agent \(B\)'s boundary, every other agent in the network must
\textbf{update its threat model}. This section formalizes the cascading
cost.

\subsection{The Network Friction
Coefficient}\label{the-network-friction-coefficient}

\begin{definition}[Network Friction]
\label{def-network-friction}

Consider a network of $N$ agents engaged in cooperative resource exchange. The \textbf{network friction coefficient} $\phi \geq 0$ parameterizes the per-transaction overhead cost due to monitoring, verification, and precautionary defense:

$$C_{\text{overhead}}(\phi) = \phi \cdot E_{\text{transaction}}$$

where $E_{\text{transaction}}$ is the energy value of a cooperative transaction. At $\phi = 0$, transactions are frictionless (perfect trust). At $\phi > 0$, each interaction bears an overhead tax.
\end{definition}

\textbf{Examples of friction overhead:} - \textbf{Biological:} Immune
surveillance costs (metabolic energy spent monitoring for pathogens,
even when healthy). - \textbf{Economic:} Transaction costs (legal fees,
contract enforcement, insurance, escrow). - \textbf{Societal:} Policing,
auditing, intelligence gathering, defensive military posture.

\subsection{Trust as a State Variable}\label{trust-as-a-state-variable}

Define the \textbf{network trust level} \(T \in [0, 1]\) as the
reciprocal transform of friction (equivalently, the complement of the
normalized friction \(\phi/(1+\phi)\)):

\[T = \frac{1}{1 + \phi}\]

\begin{itemize}
\tightlist
\item
  \(T = 1\): perfect trust (\(\phi = 0\)), frictionless cooperation.
\item
  \(T \to 0\): zero trust (\(\phi \to \infty\)), cooperative
  transactions become prohibitively expensive.
\end{itemize}

Equivalently: \(\phi = \frac{1 - T}{T}\).

\subsection{Violation-Induced Friction
Increase}\label{violation-induced-friction-increase}

Each boundary violation increases the network friction coefficient:

\begin{definition}[Friction Injection]
\label{def-friction-injection}

When a boundary violation of severity $v > 0$ (measured in energy units of damage) occurs at time $t$, the network friction coefficient updates as:

$$\phi_{t+1} = \phi_t + \eta \cdot v$$

where $\eta > 0$ is the \textbf{friction sensitivity}, which controls how strongly the network responds to violations. The sensitivity $\eta$ depends on:
\begin{itemize}
\item \textbf{Visibility:} Public violations have higher $\eta$ (all agents update simultaneously).
\item \textbf{Precedent:} First violations in a low-friction network have higher $\eta$ (trust is fragile).
\item \textbf{Network density:} Denser networks (more inter-agent connections) have higher $\eta$ (information propagates faster).
\end{itemize}

\end{definition}

\subsection{The Cost of Ambient
Friction}\label{the-cost-of-ambient-friction}

If the network has \(M\) cooperative transactions per period, each of
average value \(\bar{E}\), the \textbf{total friction tax} per period
is:

\begin{definition}[Friction Tax]
\label{def-friction-tax}

The per-period \textbf{friction tax} on a cooperative network with $M$ transactions of average value $\bar{E}$ and network friction coefficient $\phi$ is:

$$C_{\text{friction}} = \phi \cdot M \cdot \bar{E}$$

This is the aggregate overhead energy that the network dissipates on monitoring, verification, and defense, energy that could otherwise go to productive boundary maintenance.
\end{definition}

\subsection{Friction Recovery (Trust
Rebuilding)}\label{friction-recovery-trust-rebuilding}

Trust rebuilds slowly through consistent cooperative behavior:

\begin{definition}[Friction Decay]
\label{def-friction-decay}

In the absence of new violations, the friction coefficient decays toward zero at rate $\epsilon \in (0, 1)$ per period:

$$\phi_{t+1} = (1 - \epsilon) \cdot \phi_t$$

This exponential decay means friction halves every $t_{1/2} = \ln 2 / \ln(1/(1 - \epsilon))$ periods.
\end{definition}

\subsection{Combined Law of Motion}\label{combined-law-of-motion}

Combining Definition~\ref{def-friction-injection} and
Definition~\ref{def-friction-decay} into a single dynamic equation:

\begin{remark}[Friction Law of Motion]
\label{rmk-friction-law-of-motion}

The network friction coefficient $\phi_t$ evolves according to a first-order stochastic difference equation:

$$\phi_{t+1} = (1 - \epsilon)\,\phi_t + \eta\,v_t \cdot \mathds{1}_{\{\text{violation at } t\}}$$

where $\epsilon \in (0,1)$ is the trust recovery rate (Definition~\ref{def-friction-decay}), $\eta > 0$ is the friction sensitivity, $v_t > 0$ is the violation severity, and $\mathds{1}_{\{\text{violation at } t\}}$ is the indicator of a boundary violation at time $t$. In the absence of violations, friction decays geometrically toward zero. Each violation injects an additive shock $\eta v_t$ that raises the friction floor for all subsequent periods. This asymmetry (instantaneous injection, gradual recovery) is thermodynamically natural: increasing entropy is easy; decreasing it requires sustained work.
\end{remark}

\textbf{Key asymmetry:} Friction injection
(Definition~\ref{def-friction-injection}) is \textbf{instantaneous and
additive} (a single violation immediately increases \(\phi\)). Friction
recovery is \textbf{gradual and multiplicative} (rebuilding trust
requires persistent cooperative behavior over many periods). This
asymmetry is thermodynamically natural: increasing entropy (breaking
trust) is easy; decreasing entropy (rebuilding trust) requires sustained
work.

\subsection{The Cascading Cost
Theorem}\label{the-cascading-cost-theorem}

\begin{theorem}[Cascading Friction]
\label{thm-cascading-friction}

A single boundary violation of severity $v$ in a network with $M$ transactions of average value $\bar{E}$ and friction sensitivity $\eta$ produces a \textbf{total cascading cost} of:

$$C_{\text{cascade}} = \frac{\eta \cdot v \cdot M \cdot \bar{E}}{\epsilon}$$

over the full recovery period $[t_{\text{violation}}, \infty)$. The ratio of cascading cost to direct conflict cost $v$ is:

$$\frac{C_{\text{cascade}}}{v} = \frac{\eta \cdot M \cdot \bar{E}}{\epsilon}$$

This ratio scales linearly with network size ($M$) and inversely with recovery rate ($\epsilon$).
\end{theorem}

\begin{proof}
The violation at time $t_0$ injects $\Delta\phi = \eta v$ into the friction coefficient. Without further violations, the excess friction at time $t_0 + k$ is:

$$\Delta\phi_k = \eta v \cdot (1 - \epsilon)^k$$

The excess friction cost at period $k$ is $\Delta\phi_k \cdot M \cdot \bar{E}$. The total cascading cost is the sum over all future periods:

$$C_{\text{cascade}} = \sum_{k=0}^{\infty} \eta v \cdot (1 - \epsilon)^k \cdot M \cdot \bar{E} = \eta v \cdot M \cdot \bar{E} \cdot \frac{1}{\epsilon}$$

using the geometric series formula $\sum_{k=0}^{\infty} q^k = 1/(1-q)$ for $q = 1 - \epsilon \in (0, 1)$.
\end{proof}

\subsection{Numerical Illustration: The Amplification
Factor}\label{numerical-illustration-the-amplification-factor}

Consider a modest-sized economic network:

{\def\LTcaptype{none} %
\begin{longtable}[]{@{}
  >{\raggedright\arraybackslash}p{(\linewidth - 6\tabcolsep) * \real{0.2500}}
  >{\raggedright\arraybackslash}p{(\linewidth - 6\tabcolsep) * \real{0.2500}}
  >{\raggedright\arraybackslash}p{(\linewidth - 6\tabcolsep) * \real{0.2500}}
  >{\raggedright\arraybackslash}p{(\linewidth - 6\tabcolsep) * \real{0.2500}}@{}}
\toprule\noalign{}
\begin{minipage}[b]{\linewidth}\raggedright
Parameter
\end{minipage} & \begin{minipage}[b]{\linewidth}\raggedright
Symbol
\end{minipage} & \begin{minipage}[b]{\linewidth}\raggedright
Value
\end{minipage} & \begin{minipage}[b]{\linewidth}\raggedright
Interpretation
\end{minipage} \\
\midrule\noalign{}
\endhead
\bottomrule\noalign{}
\endlastfoot
Network transactions / period & \(M\) & 1,000 & A small community \\
Average transaction value & \(\bar{E}\) & 10 & Energy units per
transaction \\
Friction sensitivity & \(\eta\) & 0.01 & Friction increment per unit
violation energy (units: inverse energy) \\
Recovery rate & \(\epsilon\) & 0.05 & \textasciitilde14 periods to halve
friction \\
Violation severity & \(v\) & 50 & A significant boundary breach \\
\end{longtable}
}

\textbf{Direct cost of the violation:} \(v = 50\) energy units.

\textbf{Cascading cost:}

\[C_{\text{cascade}} = \frac{0.01 \times 50 \times 1000 \times 10}{0.05} = \frac{5000}{0.05} = 100{,}000 \text{ energy units}\]

\textbf{Amplification factor:}
\(C_{\text{cascade}} / v = 100{,}000 / 50 = 2{,}000\times\).

Scaling the same parameters to a larger ecosystem (\(M = 10{,}000\)
transactions per period, holding \(\bar{E}\), \(\eta\), \(\epsilon\),
and \(v\) fixed) yields:

\[C_{\text{cascade}} = \frac{0.01 \times 50 \times 10{,}000 \times 10}{0.05} = 1{,}000{,}000 \text{ energy units},\]

with amplification factor \(20{,}000\times\). The amplification factor
scales linearly in \(M\) and inversely in \(\epsilon\), so modest
network parameters span a band from a few times \(10^3\) to a few times
\(10^4\) across small communities and larger ecosystems.

The network-level damage from a single violation is therefore roughly
three to four orders of magnitude greater than the direct damage. This
formalizes the intuition that small lies or broken promises do not
destroy the network immediately, but they increase the ambient ``heat''
(distrust).

\subsection{Multiple Violations: The Friction
Ratchet}\label{multiple-violations-the-friction-ratchet}

When violations are recurring (rate \(\nu\) violations per period, each
of severity \(v\)), the steady-state friction coefficient is:

\[\phi_{\infty} = \frac{\nu \cdot \eta \cdot v}{\epsilon}\]

and the steady-state friction tax is:

\[C_{\text{friction}}^{\infty} = \phi_{\infty} \cdot M \cdot \bar{E} = \frac{\nu \cdot \eta \cdot v \cdot M \cdot \bar{E}}{\epsilon}\]

\begin{corollary}[Friction Ratchet]
\label{cor-friction-ratchet}

If the violation rate $\nu$ exceeds the critical threshold:

$$\nu^* = \frac{\epsilon}{\eta \cdot v} \cdot \frac{E_{\text{coop net}}}{M \cdot \bar{E}}$$

where $E_{\text{coop net}}$ is the total cooperative surplus of the network, then the friction tax exceeds the cooperative surplus and the network collapses: cooperation becomes energetically unprofitable.
\end{corollary}

This is the mathematical formalization of societal collapse: when
violations are too frequent and trust recovery is too slow, the overhead
cost of maintaining any cooperative structure exceeds its benefit, and
agents are better off in isolation or smaller, higher-trust
sub-networks.

\section{The Interaction Cost Function: Formal
Definition}\label{the-interaction-cost-function-formal-definition}

Synthesizing the results above into a single mathematical object:

\begin{definition}[Interaction Cost Function]
\label{def-interaction-cost-function}

The \textbf{interaction cost function} $\mathcal{C} : \mathbb{R}_{\geq 0}^N \times [0, \infty) \to \mathbb{R}_{\geq 0}$ maps a strategy profile and network friction state to the total energy dissipated by non-cooperative interactions:

$$\mathcal{C}(\mathbf{e}, \phi) = \underbrace{\sum_{i=1}^{N} e_i}_{\text{contest dissipation}} + \underbrace{\kappa \sum_{i=1}^{N} (\psi^{\text{off}} + \psi^{\text{def}}) \, e_i}_{\text{boundary repair (re-embodiment + surcharge)}} + \underbrace{\phi \cdot M \cdot \bar{E}}_{\text{ambient friction tax}}$$

where $\mathbf{e} = (e_1, \ldots, e_N)$ is the vector of contest investments.

At the contest Nash equilibrium:

$$\mathcal{C}^* = \Phi \cdot \frac{N-1}{N} \cdot E_j + \phi \cdot M \cdot \bar{E}$$
\end{definition}

This defines thermodynamic friction as a precise mathematical object:
the total energy dissipated by non-cooperative interactions in a
multi-agent system, decomposed into contest dissipation, boundary repair
(which re-supplies the destroyed embodied energy plus its irreversible
surcharge), and ambient friction tax.

\section{Worked Examples}\label{worked-examples}

\subsection{Example 1: Two-Agent Resource Contest
(Baseline)}\label{example-1-two-agent-resource-contest-baseline}

\textbf{Setup:} Two symmetric agents (\(\sigma_A = \sigma_B = 1\)), one
resource worth \(E_j = 100\) energy units. Boundary parameters:
\(\psi^{\text{off}} = 0.15\), \(\psi^{\text{def}} = 0.25\) (so
\(\psi = 0.4\)), \(\kappa = 2\).

\textbf{Friction multiplier:}

\[\Phi = 1 + 2 \times (0.15 + 0.25) = 1 + 2 \times 0.4 = 1.8\]

\textbf{Contest equilibrium:}

\[e_A^* = e_B^* = \frac{100}{4} = 25\]

\textbf{Total dissipation:} \(D_2 = 50\)

\textbf{Boundary damage per agent:}
\(\Delta B_i = \psi \times 25 = 0.4 \times 25 = 10\)

\textbf{Repair cost per agent:}
\(W_{\text{repair},i} = \kappa \Delta B_i = 2 \times 10 = 20\) (this
re-supplies the destroyed \(\Delta B_i = 10\) and adds a \(10\)
irreversible surcharge; \(\Delta B_i\) is not charged again)

\textbf{Total system loss:} \(2 \times (25 + 20) = 90\) (i.e.,
\(\Phi \times D_2 = 1.8 \times 50 = 90\))

\textbf{Net system value:} \(100 - 90 = +10\) energy units.

\textbf{Verdict: Net-positive, but wasteful.} At the mild baseline the
contest dissipates 90\% of the resource value; the system retains only
\(10\) of the \(100\) at stake, versus \(100\) under a cooperative split
with zero conflict cost. The two-agent contest tips net-negative once
damage becomes destructive (\(\kappa\psi > 1\), i.e.~\(\psi > 1/2\) at
\(\kappa = 2\)) or once a third contestant enters (Example 2).

\subsection{Example 2: Five-Agent
Scramble}\label{example-2-five-agent-scramble}

Same parameters as above (\(\psi = 0.4\), \(\kappa = 2\),
\(\Phi = 1.8\)), but \(N = 5\).

\[e_i^* = \frac{4 \times 100}{25} = 16, \qquad D_5 = 5 \times 16 = 80\]

\textbf{Total system loss:} \(\Phi \times 80 = 1.8 \times 80 = 144\)

\textbf{Net system value:} \(100 - 144 = -44\) energy units.

\textbf{Verdict: Net-negative.} With five contestants the threshold
falls to \(\Phi > 5/4\), which the baseline \(\Phi = 1.8\) clears
comfortably, and the system dissipates 44\% more energy than the
resource is worth. Each agent's expected contest payoff: \(100/5 = 20\)
gross from winning, minus \(16\) investment, leaving \(4 = E_j/N^2\),
before boundary repair costs.

\subsection{Example 3: The ``Profitable Bully'' (When Unilateral
Defection Pays
Temporarily)}\label{example-3-the-profitable-bully-when-unilateral-defection-pays-temporarily}

Not all defection is immediately net-negative for the defector. Consider
agents with valuations \(V_A = V_B = 100\) but asymmetric effectiveness
(\(\sigma_A = 3\), \(\sigma_B = 1\)) in the standard Tullock model:

Using the symmetric-valuation Tullock NE with effectiveness weights, the
investments are:

\[e_A^* = e_B^* = \frac{\sigma_A \sigma_B}{(\sigma_A + \sigma_B)^2} E_j = \frac{3}{16} \times 100 = 18.75\]

Winning probabilities:
\(p_A = \frac{\sigma_A}{\sigma_A + \sigma_B} = \frac{3}{4}\),
\(p_B = \frac{1}{4}\).

Agent \(A\)'s expected payoff (before boundary costs):
\(0.75 \times 100 - 18.75 = 56.25\).

Agent \(B\)'s expected payoff: \(0.25 \times 100 - 18.75 = 6.25\).

Total contest dissipation: \(2 \times 18.75 = 37.5\).

With boundary costs (\(\Phi = 1.8\)): total system loss =
\(1.8 \times 37.5 = 67.5\).

\textbf{System net:} \(100 - 67.5 = 32.5\), net-positive.

\textbf{But the bully profits disproportionately at the victim's
expense.} Even where the two-agent contest is net-positive at the system
level, the contest redistributes value sharply toward the stronger
agent: the bully nets \(56.25\) (before boundary costs) while the victim
nets only \(6.25\). The private incentive to bully therefore persists
regardless of the sign of the system net --- and it survives unchanged
into the regimes where the systemic outcome does turn net-negative
(destructive damage, \(\kappa\psi > 1\), or additional contestants).
\textbf{This is precisely why boundary constraints are necessary}
(Theorem 5 \citep{TC-I}): the bully's profit is private and short-term;
the system's loss is public, and cascading friction
(Theorem~\ref{thm-cascading-friction}) means the network-level damage
vastly exceeds the direct cost.

\section{\texorpdfstring{Decisive Contests (\(r > 1\)): Escalation
Dynamics}{Decisive Contests (r \textgreater{} 1): Escalation Dynamics}}\label{decisive-contests-escalation-dynamics}

\subsection{The Impact of
Decisiveness}\label{the-impact-of-decisiveness}

When \(r > 1\), the contest becomes increasingly ``all-or-nothing'': a
marginal increase in investment translates into a larger increase in
winning probability, and agents invest more aggressively as a result.

For symmetric agents in the moderate-decisiveness regime
\(r \in (0, N/(N-1)]\), the interior first-order condition yields the
unique symmetric pure-strategy Nash equilibrium

\[e_i^* = \frac{r(N-1)}{N^2} \cdot E_j, \qquad D_N = \frac{r(N-1)}{N} \cdot E_j.\]

At the upper boundary \(r = N/(N-1)\), dissipation reaches the full
resource value (\(D_N = E_j\)) and per-agent expected payoff falls to
zero. For \(r > N/(N-1)\), the interior critical point is no longer a
maximum of agent \(i\)'s payoff: the symmetric pure-strategy Nash
equilibrium of the unconstrained Tullock contest ceases to exist
\citep{BayeKovenock-deVries1994}, and the equilibrium of the unbounded
game shifts into mixed strategies in which expected total dissipation
equals \(E_j\) (full rent dissipation).

\subsection{Physical Budget
Constraints}\label{physical-budget-constraints}

The unbounded contest is a useful theoretical limit, but it does not
describe the systems this paper considers. In our framework every agent
is an open thermodynamic system with finite metabolic supply: contest
investment \(e_i\) is real energy diverted from boundary maintenance and
productive use, and is bounded above by an agent-specific physical
budget \(\bar e_i > 0\) (a function of metabolic capacity, stored
reserves, and boundary integrity \(B_i\)). The contest action space is
therefore the compact interval \([0, \bar e_i]\) rather than
\([0, \infty)\).

In the moderate regime \(r \leq N/(N-1)\), the interior NE
\(e_i^* = r(N-1) E_j/N^2\) remains the equilibrium whenever it lies
inside the budget set (\(e_i^* \leq \bar e_i\)); otherwise the
equilibrium is at the cap, examined below.

In the decisive regime \(r > N/(N-1)\), the interior FOC has no maximum
on the unbounded action space, but the budget constraint compactifies
the strategy space and restores existence of a pure-strategy
equilibrium, located at the budget cap.

\begin{proposition}[Budget-Constrained Decisive Contest]
\label{prop-budget-constrained-decisive}

Consider $N \geq 2$ symmetric agents ($\sigma_i = 1$) with common physical budget $\bar e \in (0, E_j/N]$ contesting a resource of value $E_j$ under decisiveness $r > N/(N-1)$. The symmetric profile $e_i = \bar e$ for all $i$ is a pure-strategy Nash equilibrium of the budget-constrained contest. Total contest dissipation is

$$D_N = N \bar e \leq E_j,$$

with equality when $\bar e = E_j/N$. Each agent's expected payoff is $\Pi_i^* = E_j/N - \bar e \geq 0$.
\end{proposition}

\begin{proof}
At the symmetric profile $e_i = \bar e$, agent $i$'s success probability is $1/N$ and expected payoff is $E_j/N - \bar e \geq 0$. Two classes of unilateral deviation must be checked.

\emph{Upward deviations} are infeasible: the budget cap $\bar e$ binds.

\emph{Downward deviations} to $e_i' \in [0, \bar e)$. The marginal payoff at $e_i = \bar e$ when others play $\bar e$ is $E_j r(N-1)/(N^2 \bar e) - 1$. Because $\bar e \leq E_j/N$ and $r > N/(N-1)$,

$$\frac{E_j \, r(N-1)}{N^2 \, \bar e} \;\geq\; \frac{r(N-1)}{N} \;>\; 1,$$

so the marginal payoff is strictly positive and the payoff is locally increasing as $e_i \to \bar e^-$; small downward deviations strictly reduce payoff. For the global comparison over all $e_i' \in [0, \bar e]$, write the deviation payoff as $g(e) = \frac{e^r}{e^r + (N-1)\bar e^r} E_j - e$. Its derivative is $g'(e) = h(e) - 1$ with $h(e) = \frac{E_j \, r (N-1) \bar e^r e^{r-1}}{(e^r + (N-1)\bar e^r)^2}$. For $r > 1$, $h(0) = 0$ and $h$ is single-peaked in $e$; since $h(\bar e) = \frac{E_j \, r(N-1)}{N^2 \bar e} > 1$ (shown above), $h$ crosses $1$ at most once on $[0, \bar e]$, and from below. Hence $g$ is decreasing then increasing on $[0, \bar e]$, so its maximum over the interval is attained at an endpoint. Comparing them, $g(\bar e) = E_j/N - \bar e \geq 0 = g(0)$, so the cap is the global best response. Hence $e_i = \bar e$ for all $i$ is a pure-strategy Nash equilibrium.
\end{proof}

The interior NE of the moderate regime and the corner NE of the decisive
regime meet continuously at \(r = N/(N-1)\) when \(\bar e = E_j/N\): the
interior formula gives
\(e_i^* = (N-1) E_j/N^2 \cdot N/(N-1) = E_j/N = \bar e\).

\subsection{Net-Negativity Under Budget
Constraints}\label{net-negativity-under-budget-constraints}

The friction multiplier \(\Phi\) applies identically in either regime,
so the system loss is

\[L_{\text{system}}^{(N)} = \Phi \cdot D_N,\]

with \(D_N = r(N-1) E_j/N\) in the moderate regime and
\(D_N = N \bar e\) in the budget-bound decisive regime.

\begin{corollary}[Escalation under Physical Budgets]
\label{cor-escalation-instability}

For $N \geq 2$ symmetric agents with budget cap $\bar e$:

(a) \emph{Moderate regime} ($r \leq N/(N-1)$, interior NE feasible): conflict is net-negative whenever $\Phi r (N-1)/N > 1$. In particular, $r = 1$ recovers Theorem~\ref{thm-net-negative-conflict}.

(b) \emph{Decisive regime} ($r > N/(N-1)$, equilibrium at cap): conflict is net-negative whenever $\Phi N \bar e > E_j$, i.e., whenever the cap exceeds $E_j/(\Phi N)$. With the baseline $\Phi = 1.8$ and $N = 2$, this requires $\bar e > E_j/3.6$, a budget level that is small relative to the resource value but routinely exceeded by adversaries committing significant fractions of their reserves to the contest.

In both regimes, sufficiently decisive contests with non-negligible boundary damage produce dissipation exceeding the contested resource value before the contest is even resolved.
\end{corollary}

\textbf{Physical interpretation.} Arms races and high-stakes
confrontations are decisive contests with \(r > 1\) and budget
commitments approaching \(\bar e\). The decisive-regime result captures
what unbounded Tullock cannot: agents do not bid to mathematical
infinity, they exhaust their physical reserves. The cap is set by
metabolism, treasury, or boundary integrity, not by strategic
moderation. Once \(\Phi\) multiplies the cap-bound dissipation, the
total energetic cost of the contest exceeds the value of the contested
resource, recovering the arms-race conclusion that history confirms.

\section{Discussion}\label{discussion}

\subsection{Physical Interpretation}\label{physical-interpretation}

The results establish a complete energetic cost structure for
adversarial resource competition. The logical chain proceeds as follows:
resource collision (established by \emph{Necessary Constraints}
\citep{TC-I}) forces a contest; the contest dissipates energy
proportional to the resource value
(Theorem~\ref{thm-dissipation-scaling}); boundary damage and
thermodynamically irreversible repair amplify the cost through the
friction multiplier (Theorem~\ref{thm-net-negative-conflict}); and
network cascading propagates a single violation across all future
cooperative transactions (Theorem~\ref{thm-cascading-friction}).

The net-negative theorem (Theorem~\ref{thm-net-negative-conflict}) is
the central result: the total energy dissipated by conflict (contest
effort, boundary damage, and thermodynamically irreversible repair)
exceeds the value of the contested resource whenever \(\Phi > N/(N-1)\).
This threshold is crossed either by destructive boundary damage in the
two-agent case (\(\Phi > 2 \Leftrightarrow \kappa\psi > 1\)) or, for any
positive damage, by three or more contestants
(\(\Phi > N/(N-1) \leq 3/2\)). In these regimes conflict is not merely
suboptimal; it is thermodynamically wasteful at a system level: the
system would be better off if the resource simply did not exist, because
the act of fighting over it destroys more value than the resource
contains. Even below the threshold --- the mild two-agent baseline
dissipates 90\% of the resource value --- conflict is severely
value-eroding relative to the cooperative split.

The network cascading result (Theorem~\ref{thm-cascading-friction})
extends this from bilateral to multilateral settings. A single violation
does not merely cost the immediate parties; it injects friction into the
cooperative network that degrades all subsequent interactions, with
illustrative amplification ratios of roughly \(10^3\) to \(10^4\) under
modest network parameters. This makes the true cost of defection far
larger than the immediate bilateral loss. The asymmetry between friction
injection (instantaneous, additive) and friction recovery (gradual,
multiplicative) is thermodynamically natural: increasing entropy is
easy; decreasing it requires sustained work.

Together, these results provide the physical foundation for the
energy-denominated payoff matrix used in \emph{Cooperative Equilibrium}
\citep{TC-V}. The costly mutual-defection payoff in that matrix is not
assumed; it is a consequence of the friction analysis proved here, and
it becomes strictly negative (\(P < 0\)) in exactly the regime this
paper identifies --- destructive boundary damage (\(\kappa\psi > 1\)) or
three or more contestants.

The theoretical amplification ratios obtained in the worked illustration
(roughly \(10^3\) to \(10^4\)) exceed those reported in empirical
studies of specific conflict events (e.g., the approximately 4:1 supply
chain amplification measured in the 2017 NotPetya cyberattack
\citep{Crosignani2023}). This discrepancy is expected and systematic,
for three reasons. First, the model computes costs for contests played
to equilibrium; empirical conflicts typically terminate early (through
settlement, withdrawal, or capitulation) once the accumulating costs
become apparent, so empirical data captures the cost of partially played
contests. Second, the model is denominated in energy, not money;
monetary estimates omit energy flows that have no market price
(metabolic costs of stress and displacement, thermodynamic cost of
rebuilding to lower specification) and undercount flows that are
nominally priced but practically unattributable due to the
interdependency of causal factors (reduced productivity under elevated
uncertainty, opportunity costs of diverted capital, sustained cognitive
load across affected populations). Third, empirical observation windows
are finite, while the cascading friction integral extends over the full
recovery period, including the long tail of elevated ambient friction
that no finite-horizon study captures. Each of these biases operates in
the same direction: empirical dollar-denominated estimates are
systematic lower bounds on the total energy-denominated cost that the
theoretical model computes.

Moreover, even the friction model itself is conservative. The cost
structure derived here (contest expenditure, boundary damage, repair,
and network cascading) accounts for the energy dissipated by the
conflict process, but does not account for the \emph{informational}
value of what conflict destroys. \emph{Accumulated Negentropy}
\citep{TC-III} proves that complex dissipative structures accumulate
organizational information over time
(\(\mathcal{N}(T) = \int_0^T \dot{W}_{\text{order}}(t)\,dt\)) that
cannot be reconstituted by re-spending the original energy investment.
When conflict destroys such structures, the true loss includes not only
the repair cost \(\kappa\) captured by the friction multiplier but also
the accumulated negentropy that no repair process recovers. The friction
analysis proved here is therefore a lower bound on the total
thermodynamic cost of conflict; incorporating accumulated negentropy
would increase the effective cost by additional orders of magnitude.

\subsection{Testable Predictions}\label{testable-predictions}

The model generates falsifiable predictions with explicit failure
conditions:

{\def\LTcaptype{none} %
\begin{longtable}[]{@{}
  >{\raggedright\arraybackslash}p{(\linewidth - 4\tabcolsep) * \real{0.3333}}
  >{\raggedright\arraybackslash}p{(\linewidth - 4\tabcolsep) * \real{0.3333}}
  >{\raggedright\arraybackslash}p{(\linewidth - 4\tabcolsep) * \real{0.3333}}@{}}
\toprule\noalign{}
\begin{minipage}[b]{\linewidth}\raggedright
Prediction
\end{minipage} & \begin{minipage}[b]{\linewidth}\raggedright
Source
\end{minipage} & \begin{minipage}[b]{\linewidth}\raggedright
Falsified If\ldots{}
\end{minipage} \\
\midrule\noalign{}
\endhead
\bottomrule\noalign{}
\endlastfoot
Contest dissipation approaches the resource value as contestant number
grows: \(D_N/E_j = (N-1)/N \to 1\) &
Theorem~\ref{thm-dissipation-scaling} & Empirical resource contests with
\(N \geq 10\) (e.g., multi-bidder rent-seeking competitions)
consistently dissipate less than \(80\%\) of the contested value, where
the model predicts at least \(90\%\) \\
Conflict is net-negative (\(\Phi > N/(N-1)\)): two-agent contests once
\(\kappa\psi > 1\) (e.g.~\(\kappa = 2\),
\(\psi = \psi^{\text{off}} + \psi^{\text{def}} \geq 0.5\)), and every
contest with \(N \geq 3\) under any positive boundary damage &
Theorem~\ref{thm-net-negative-conflict} & Empirical data on conflict
costs in resource competition (military expenditure vs.~resource value,
litigation costs vs.~settlement values) consistently show total costs
below 100\% of the contested resource value when boundary damage is
non-negligible and contestants are numerous \\
Network amplification of violation costs at ratios of roughly \(10^3\)
to \(10^4\) & Theorem~\ref{thm-cascading-friction} & Supply chain
disruption data, financial contagion events, or ecological cascade
measurements show amplification ratios consistently below \(10^2\) in
densely connected cooperative networks \\
Cooperation strictly dominates conflict for all admissible friction
parameters (\(\Phi \geq 1\)) whenever resource collision occurs &
Theorem~\ref{thm-cooperation-dominance} & Experimental resource-sharing
games with calibrated friction parameters consistently find conflict
outcomes that Pareto-dominate cooperative outcomes \\
Decisive contests (\(r > N/(N-1)\)) under finite physical budgets
dissipate the full budget commitment \(N\bar e\), exceeding \(E_j\) once
amplified by \(\Phi\) whenever \(\bar e > E_j/(\Phi N)\) &
Corollary~\ref{cor-escalation-instability} & Arms-race expenditure data
shows total military spending below the strategic value of the contested
territory or resource \\
\end{longtable}
}

The strongest falsification of the model would be a class of physically
realistic conflict scenarios, with non-negligible boundary damage and
thermodynamically irreversible repair, in which total conflict costs are
consistently below the contested resource value. The model's
architecture makes this structurally difficult (\(\Phi > 1\) whenever
boundary damage is positive), but a counterexample would require
revision of the damage model.

\subsection{Limitations}\label{limitations}

The contest model uses the Tullock success function, which assumes
agents compete by expending effort that enters a ratio-form contest
success function. Real conflicts may involve threshold effects, arms
races with discontinuous outcomes, or asymmetric information about
opponent capabilities. The boundary damage model assumes linear damage
fractions (\(\psi^{\text{off}}, \psi^{\text{def}}\) are constants);
nonlinear or state-dependent damage functions would refine the friction
multiplier estimates but would not eliminate the fundamental result that
repair costs exceed the damage they restore (which follows from the
Second Law, not from linearity).

The network cascading analysis assumes a homogeneous network topology
with uniform friction transmission. Heterogeneous networks with variable
connectivity would yield a distribution of amplification ratios rather
than a single estimate. The friction sensitivity parameter \(\eta\) and
recovery rate \(\epsilon\) are treated as constants; in practice, both
may depend on the history of violations, network structure, and agent
characteristics.

The repair multiplier \(\kappa\) is bounded below by 1 (thermodynamic
irreversibility) but its precise value is application-specific. The
worked examples use \(\kappa = 2\) as a deliberately conservative upper
baseline: read as an \emph{energy ratio} (metabolic energy to regenerate
structure ÷ embodied energy of the structure restored), the biological
synthesis evidence cleanly supports a range of roughly 1.2--1.7, and
\(\kappa = 2\) sits \emph{above} that range, folding in the integrated
inflammatory overhead of a full healing episode rather than resting on
the synthesis ratio alone. Because the two-agent net-negativity verdict
turns on \(\kappa\psi\) relative to 1, we report its sensitivity around
the sign-flip line explicitly; the \(N\)-agent result and the empirical
cost-to-value ratios carry the qualitative conclusion independent of the
precise value of \(\kappa\). Other domains (institutional trust,
financial system repair) may have substantially higher values.

\textbf{Scope.} The model applies to any system of agents that maintain
boundaries and compete for finite resources, with no restriction on
species, substrate, or scale. The results address the energy cost of
conflict, not the mechanism by which cooperation is established or
enforced. Subsequent papers in this series address enforcement through
game-theoretic equilibrium selection (\emph{Cooperative Equilibrium}
\citep{TC-V}) and the dynamics of cooperative stability (\emph{Value
Dynamics} \citep{TC-VI}).

\subsection{Applications}\label{applications}

The model applies wherever self-maintaining systems engage in
adversarial resource competition:

\begin{enumerate}
\def\labelenumi{\arabic{enumi}.}
\item
  \textbf{Military conflict and territorial disputes.} The friction
  multiplier \(\Phi\) provides a computable estimate of total conflict
  cost relative to the contested resource value. The net-negative
  theorem predicts that wars over territory or resources dissipate more
  energy (economic output, infrastructure, human capital) than the
  contested resources are worth, a prediction consistent with empirical
  analyses of conflict costs \citep{GarfinkelSkaperdas2012}.
  \citet{StiglitzBilmes2008} estimate the total economic cost of the
  Iraq War at approximately \$3 trillion (including long-term veteran
  care, debt service, and macroeconomic disruption); against annual
  Iraqi oil revenues of approximately \$96 billion \citep{IMF2024Iraq},
  this yields a cost-to-annual-resource-value ratio exceeding 30:1. The
  escalation result (Corollary~\ref{cor-escalation-instability})
  predicts that decisive contests under finite physical budgets exhaust
  the budgeted energy at a rate that, once amplified by \(\Phi\),
  exceeds the strategic value of the contested resource.
\item
  \textbf{Litigation and dispute resolution.} Legal contests are Tullock
  contests with \(r\) determined by the quality of evidence and legal
  representation. The model predicts that total litigation costs
  (attorney fees, court costs, productivity loss, relationship damage)
  should exceed settlement values under adversarial proceedings with
  non-negligible boundary damage. \citet{HannafordAgor2013} report that
  median litigation costs range from \$43,000 to \$122,000, while
  \citet{HannafordAgor2015} found that 75\% of civil judgments are below
  \$5,200, yielding cost-to-value ratios exceeding 8:1 for the majority
  of cases. These findings are consistent with the net-negative
  prediction: the contest expenditure alone consumes the resource value,
  before accounting for the additional friction of relationship damage
  and productivity loss.
\item
  \textbf{Cybersecurity and network defense.} Network intrusion is a
  resource contest in which attackers expend resources to breach system
  boundaries (firewalls, access controls, encryption) while defenders
  allocate resources to detection and hardening. The boundary integrity
  parameter \(B_i\) maps to system hardening investment, and breach
  repair (incident response, system restoration, credential rotation)
  consistently costs multiples of the bypassed security investment,
  confirming \(\kappa > 1\) in this domain. The network cascading result
  (Theorem~\ref{thm-cascading-friction}) predicts that a single
  compromised node in a supply chain propagates friction across all
  dependent systems. \citet{Crosignani2023} document this effect in the
  2017 NotPetya attack: disruptions at directly compromised firms
  propagated through supply chains to their customers, with aggregate
  downstream losses approximately four times larger than the direct
  losses to the initially targeted firms.
\item
  \textbf{Ecological competition and niche partitioning.} Interspecific
  competition for overlapping resource niches is a biological Tullock
  contest where the ``boundary'' is the organism's physiological
  integrity and territorial defense. The cooperation dominance result
  predicts that niche partitioning (the ecological analog of boundary
  constraints) should yield higher total biomass than direct
  competition. A meta-analysis of 44 independent experiments found that
  species mixtures produce on average 1.7 times more biomass than
  monocultures, with diverse polycultures outperforming the average
  monoculture in 79\% of experiments \citep{Cardinale2007}. This is
  consistent with the model's prediction: partitioning resources (the
  ecological analog of mutual boundary constraints) reduces contest
  dissipation and increases the total energy captured by the community.
\end{enumerate}

\subsection{Future Work}\label{future-work}

\textbf{Nonlinear damage models.} The linear damage fractions
\(\psi^{\text{off}}\) and \(\psi^{\text{def}}\) could be replaced with
state-dependent functions \(\psi^{\text{off}}(B_i, e_i)\) and
\(\psi^{\text{def}}(B_j, e_i)\) that capture threshold effects (a
boundary remains intact until a critical stress is exceeded, then fails
catastrophically) or hardening (repeated conflict strengthens the
boundary). Such extensions would refine the friction multiplier
landscape without altering the qualitative net-negative result.

\textbf{Heterogeneous network topology.} The cascading friction analysis
assumes a homogeneous network. Extension to scale-free, small-world, or
hierarchical network topologies would yield topology-dependent
amplification distributions, potentially identifying network structures
that are more or less resilient to violation cascades. Embedding the
cascade in a production network, where agents' outputs feed other
agents' inputs, would make amplification a function of the dependency
graph: friction injected at hub nodes (high out-degree) would amplify
far more than friction injected at leaf nodes. Predicting which network
positions are most destructive under defection requires the spectral
properties of the aggregated input-output matrix; such an extension
would produce structured predictions about which agent defections pose
systemic risk.

\textbf{Stochastic conflict dynamics.} The current model is
deterministic. A stochastic extension incorporating random contest
outcomes, imperfect information about opponent capabilities, and
uncertain boundary-damage magnitudes would connect the model to the
literatures on stochastic games and Bayesian conflict models.

\textbf{Empirical calibration.} The damage coefficients
\(\psi^{\text{off}}\), \(\psi^{\text{def}}\), the repair multiplier
\(\kappa\), and the network parameters \(\eta\), \(\epsilon\) can be
estimated from empirical data across domains: biological wound-healing
costs, military expenditure-to-GDP ratios in conflict zones, litigation
cost data, and financial contagion propagation rates. Such calibration
would convert the qualitative predictions into domain-specific
quantitative forecasts.

\section{Conclusion}\label{conclusion}

We have quantified the energy cost of adversarial resource competition
and proved that unconstrained conflict is a net-negative thermodynamic
outcome across a broad and physically realistic regime. The
net-negativity condition \(\Phi > N/(N-1)\)
(Theorem~\ref{thm-net-negative-conflict}) guarantees that the total
energy dissipated by conflict (contest expenditure, boundary damage, and
irreversible repair) exceeds the value of the contested resource ---
reached in the two-agent case by destructive damage (\(\kappa\psi > 1\))
and, for any positive damage, by three or more contestants. Contest
dissipation alone approaches the full resource value as the number of
contestants grows (Theorem~\ref{thm-dissipation-scaling}), and
cooperation strictly dominates conflict for all admissible parameters
(Theorem~\ref{thm-cooperation-dominance}). Network cascading amplifies
the cost of a single violation by an amount that scales with network
size and inversely with trust recovery rate, illustratively roughly
three to four orders of magnitude under modest network parameters
(Theorem~\ref{thm-cascading-friction}), formalizing the thermodynamic
asymmetry between trust destruction and trust recovery.

These results provide the energy-denominated cost structure for the
subsequent papers in this series. The costly mutual-defection payoff in
\emph{Cooperative Equilibrium} \citep{TC-V} --- strictly negative
(\(P < 0\)) once damage is destructive or contestants are numerous ---
is a direct consequence of the friction analysis proved here, and the
interaction cost function
(Definition~\ref{def-interaction-cost-function}), which synthesizes
contest dissipation, boundary damage, repair costs, and ambient
friction, enters both the payoff matrices of \emph{Cooperative
Equilibrium} and the dynamical model of \emph{Value Dynamics}
\citep{TC-VI}.

\addcontentsline{toc}{section}{References}
\begin{thebibliography}{28}
\providecommand{\natexlab}[1]{#1}
\providecommand{\url}[1]{\texttt{#1}}
\expandafter\ifx\csname urlstyle\endcsname\relax
  \providecommand{\doi}[1]{doi: #1}\else
  \providecommand{\doi}{doi: \begingroup \urlstyle{rm}\Url}\fi

\bibitem[Lostracco(2026{\natexlab{a}})]{TC-I}
Keith Lostracco.
\newblock Thermodynamics of cooperation: Necessary constraints,
  2026{\natexlab{a}}.
\newblock URL \url{https://doi.org/10.5281/zenodo.19635523}.
\newblock \doi{10.5281/zenodo.19635523}.

\bibitem[Tullock(1980)]{Tullock1980}
Gordon Tullock.
\newblock Efficient rent seeking.
\newblock In James~M. Buchanan, Robert~D. Tollison, and Gordon Tullock,
  editors, \emph{Toward a Theory of the Rent-Seeking Society}, pages 97--112.
  Texas A\&M University Press, 1980.

\bibitem[Lostracco(2026{\natexlab{b}})]{TC-V}
Keith Lostracco.
\newblock Thermodynamics of cooperation: Cooperative equilibrium,
  2026{\natexlab{b}}.
\newblock URL \url{https://doi.org/10.5281/zenodo.19635856}.
\newblock \doi{10.5281/zenodo.19635856}.

\bibitem[Skaperdas(1996)]{Skaperdas1996}
Stergios Skaperdas.
\newblock Contest success functions.
\newblock \emph{Economic Theory}, 7\penalty0 (2):\penalty0 283--290, 1996.
\newblock \doi{10.1007/BF01213906}.

\bibitem[Hirshleifer(1991)]{Hirshleifer1991}
Jack Hirshleifer.
\newblock The technology of conflict as an economic activity.
\newblock \emph{American Economic Review}, 81\penalty0 (2):\penalty0 130--134,
  1991.

\bibitem[Garfinkel and Skaperdas(2012)]{GarfinkelSkaperdas2012}
Michelle~R. Garfinkel and Stergios Skaperdas, editors.
\newblock \emph{The Oxford Handbook of the Economics of Peace and Conflict}.
\newblock Oxford University Press, 2012.
\newblock \doi{10.1093/oxfordhb/9780195392777.001.0001}.
\newblock URL \url{https://sites.socsci.uci.edu/~mrgarfin/OUP/}.

\bibitem[Konrad(2009)]{Konrad2009}
Kai~A. Konrad.
\newblock \emph{Strategy and Dynamics in Contests}.
\newblock Oxford University Press, 2009.

\bibitem[Schr{\"o}dinger(1944)]{Schrodinger1944}
Erwin Schr{\"o}dinger.
\newblock \emph{What Is Life? The Physical Aspect of the Living Cell}.
\newblock Cambridge University Press, 1944.

\bibitem[Prigogine(1978)]{Prigogine1978}
Ilya Prigogine.
\newblock Time, structure, and fluctuations.
\newblock \emph{Science}, 201\penalty0 (4358):\penalty0 777--785, 1978.
\newblock \doi{10.1126/science.201.4358.777}.

\bibitem[England(2013)]{England2013}
Jeremy~L. England.
\newblock Statistical physics of self-replication.
\newblock \emph{The Journal of Chemical Physics}, 139\penalty0 (12):\penalty0
  121923, 2013.
\newblock \doi{10.1063/1.4818538}.

\bibitem[Friston(2013)]{Friston2013}
Karl Friston.
\newblock Life as we know it.
\newblock \emph{Journal of the Royal Society Interface}, 10\penalty0
  (86):\penalty0 20130475, 2013.
\newblock \doi{10.1098/rsif.2013.0475}.

\bibitem[Maturana and Varela(1980)]{Maturana1980}
Humberto~R. Maturana and Francisco~J. Varela.
\newblock \emph{Autopoiesis and Cognition: The Realization of the Living}.
\newblock D. Reidel, 1980.
\newblock Original work published 1972.

\bibitem[Jackson(2008)]{Jackson2008}
Matthew~O. Jackson.
\newblock \emph{Social and Economic Networks}.
\newblock Princeton University Press, 2008.
\newblock ISBN 978-0-691-13440-6.

\bibitem[Ostrom(1990)]{Ostrom1990}
Elinor Ostrom.
\newblock \emph{Governing the Commons: The Evolution of Institutions for
  Collective Action}.
\newblock Cambridge University Press, 1990.

\bibitem[Lotka(1922)]{Lotka1922}
Alfred~J. Lotka.
\newblock Contribution to the energetics of evolution.
\newblock \emph{Proceedings of the National Academy of Sciences}, 8\penalty0
  (6):\penalty0 147--151, 1922.
\newblock \doi{10.1073/pnas.8.6.147}.

\bibitem[Georgescu-Roegen(1971)]{GeorgescuRoegen1971}
Nicholas Georgescu-Roegen.
\newblock \emph{The Entropy Law and the Economic Process}.
\newblock Harvard University Press, 1971.

\bibitem[Georgescu-Roegen(1979)]{GeorgescuRoegen1979}
Nicholas Georgescu-Roegen.
\newblock Comments on the papers by {Daly} and {Stiglitz}.
\newblock In V.~Kerry Smith, editor, \emph{Scarcity and Growth Reconsidered},
  pages 95--105. Johns Hopkins University Press, Baltimore, 1979.

\bibitem[Reeds et~al.(1982)Reeds, Wahle, and Haggarty]{Reeds1982EnergyCosts}
P.~J. Reeds, K.~W.~J. Wahle, and P.~Haggarty.
\newblock Energy costs of protein and fatty acid synthesis.
\newblock \emph{Proceedings of the Nutrition Society}, 41\penalty0
  (2):\penalty0 155--159, 1982.
\newblock \doi{10.1079/PNS19820025}.

\bibitem[Demling(2009)]{Demling2009}
Robert~H. Demling.
\newblock Nutrition, anabolism, and the wound healing process: An overview.
\newblock \emph{Eplasty}, 9:\penalty0 e9, 2009.
\newblock PMID: 19274069; PMCID: PMC2642618.

\bibitem[Baye et~al.(1994)Baye, Kovenock, and \{de
  Vries\}]{BayeKovenock-deVries1994}
\{MR (Michael)\} Baye, \{DJ (Dan)\} Kovenock, and Casper \{de Vries\}.
\newblock The solution to the tullock rent-seeking game when r > 2:
  mixed-strategy equilibria and mean dissipation rates.
\newblock \emph{Public Choice}, pages 363--380, 1994.
\newblock ISSN 0048-5829.
\newblock \doi{10.1007/BF01053238}.

\bibitem[Crosignani et~al.(2023)Crosignani, Macchiavelli, and
  Silva]{Crosignani2023}
Matteo Crosignani, Marco Macchiavelli, and Andr{\'e}~F. Silva.
\newblock Pirates without borders: The propagation of cyberattacks through
  firms' supply chains.
\newblock \emph{Journal of Financial Economics}, 147\penalty0 (2):\penalty0
  432--448, 2023.
\newblock \doi{10.1016/j.jfineco.2022.12.002}.

\bibitem[Lostracco(2026{\natexlab{c}})]{TC-III}
Keith Lostracco.
\newblock Thermodynamics of cooperation: Accumulated negentropy,
  2026{\natexlab{c}}.
\newblock URL \url{https://doi.org/10.5281/zenodo.19635836}.
\newblock \doi{10.5281/zenodo.19635836}.

\bibitem[Lostracco(2026{\natexlab{d}})]{TC-VI}
Keith Lostracco.
\newblock Thermodynamics of cooperation: Value dynamics, 2026{\natexlab{d}}.
\newblock URL \url{https://doi.org/10.5281/zenodo.19635865}.
\newblock \doi{10.5281/zenodo.19635865}.

\bibitem[Stiglitz and Bilmes(2008)]{StiglitzBilmes2008}
Joseph~E. Stiglitz and Linda~J. Bilmes.
\newblock \emph{The Three Trillion Dollar War: The True Cost of the {Iraq}
  Conflict}.
\newblock W. W. Norton, 2008.

\bibitem[{International Monetary Fund}(2024)]{IMF2024Iraq}
{International Monetary Fund}.
\newblock Iraq: 2024 {Article IV} consultation -- staff report.
\newblock Technical Report Country Report No.\ 24/128, International Monetary
  Fund, May 2024.
\newblock URL
  \url{https://www.elibrary.imf.org/view/journals/002/2024/128/article-A001-en.xml}.

\bibitem[Hannaford-Agor and Waters(2013)]{HannafordAgor2013}
Paula Hannaford-Agor and Nicole~L. Waters.
\newblock Estimating the cost of civil litigation.
\newblock \emph{Caseload Highlights}, 20\penalty0 (1):\penalty0 1--11, 2013.

\bibitem[Hannaford-Agor et~al.(2015)Hannaford-Agor, Graves, and
  Miller]{HannafordAgor2015}
Paula Hannaford-Agor, Scott Graves, and Shelley~Spacek Miller.
\newblock Civil justice initiative: The landscape of civil litigation in state
  courts.
\newblock 2015.
\newblock URL \url{https://scholarship.law.wm.edu/facpubs/2390}.

\bibitem[Cardinale et~al.(2007)Cardinale, Wright, Cadotte, Carroll, Hector,
  Srivastava, Loreau, and Weis]{Cardinale2007}
Bradley~J. Cardinale, Justin~P. Wright, Marc~W. Cadotte, Ian~T. Carroll, Andy
  Hector, Diane~S. Srivastava, Michel Loreau, and Jerome~J. Weis.
\newblock Impacts of plant diversity on biomass production increase through
  time because of species complementarity.
\newblock \emph{Proceedings of the National Academy of Sciences}, 104\penalty0
  (46):\penalty0 18123--18128, 2007.
\newblock \doi{10.1073/pnas.0709069104}.

\end{thebibliography}
\end{document}
